
\documentclass[UTF8,zihao=-4]{ctexart}

\usepackage[a4paper,left=2cm,right=2cm,top=2cm,bottom=2cm]{geometry}
\usepackage{amsmath,amssymb}
\usepackage{booktabs}
\usepackage{array}
\usepackage{enumitem}
\usepackage{longtable}

\setlength{\parindent}{2em}
\setlength{\parskip}{0.25em}
\setlength{\abovedisplayskip}{4pt}
\setlength{\belowdisplayskip}{4pt}
\setlength{\abovedisplayshortskip}{2pt}
\setlength{\belowdisplayshortskip}{2pt}
\renewcommand{\arraystretch}{1.12}

\newcommand{\problem}[1]{\subsection*{#1}}
\newcommand{\timu}{\noindent\textbf{题目：}}
\newcommand{\jieda}{\noindent\textbf{解答：}}

\begin{document}

\begin{center}
    {\LARGE \textbf{误差理论与数据处理}}
\end{center}

\section*{第一章\quad 绪论}
\problem{例1-1}

\timu 某台标称示值范围为 $0\sim150\,\mathrm{V}$ 的电压表（即满量程为 $150\,\mathrm{V}$），在示值为 $100\,\mathrm{V}$ 处，用标准电压表检定得到的电压表实际示值为 $99.4\,\mathrm{V}$，求使用该电压表在测得示值为 $100\,\mathrm{V}$ 时的绝对误差、相对误差和引用误差？

\jieda 该电压表在 $100\,\mathrm{V}$ 处的
\begin{equation*}
\begin{aligned}
    \text{绝对误差}
    &=100\,\mathrm{V}-99.4\,\mathrm{V}
    =0.6\,\mathrm{V}\\
    \text{相对误差}
    &=
    \frac{0.6\,\mathrm{V}}{99.4\,\mathrm{V}}\times100\%
    \approx
    \frac{0.6\,\mathrm{V}}{100\,\mathrm{V}}\times100\%
    =0.6\%\\
    \text{引用误差}
    &=
    \frac{100\,\mathrm{V}-99.4\,\mathrm{V}}{150\,\mathrm{V}}\times100\%
    =0.4\%
\end{aligned}
\end{equation*}

\section*{第二章\quad 误差的基本性质与处理}

\problem{例2-1}
\timu 测量某物理量 $10$ 次，得到结果见表2-1，求算术平均值。

\begin{center}
\scriptsize
\setlength{\tabcolsep}{7pt}
\begin{tabular}{c|cccccccccc}
\toprule
序号 & 1 & 2 & 3 & 4 & 5 & 6 & 7 & 8 & 9 & 10\\
\midrule
$l_i/\mathrm{mm}$ & 1879.64 & 1879.69 & 1879.60 & 1879.69 & 1879.57 & 1879.62 & 1879.64 & 1879.65 & 1879.66 & 1879.62\\
$\Delta l_i/\mathrm{mm}$ & $-0.01$ & $+0.04$ & $-0.05$ & $+0.04$ & $-0.08$ & $-0.03$ & $-0.01$ & 0 & $+0.01$ & $-0.03$\\
$v_i/\mathrm{mm}$ & 0 & $+0.05$ & $-0.04$ & $+0.05$ & $-0.07$ & $-0.02$ & 0 & $+0.01$ & $+0.02$ & $-0.02$\\
\bottomrule
\end{tabular}
\end{center}

\jieda 任选参考值 $l_0=1879.65\,\mathrm{mm}$，计算差值 $\Delta l_i$ 和 $\Delta\bar{x}_0$，有
\begin{equation*}
    \Delta\bar{x}_0=\frac{\sum_{i=1}^{10}\Delta l_i}{10}=-0.01\,\mathrm{mm}
\end{equation*}
所以
\begin{equation*}
    \bar{x}=l_0+\Delta\bar{x}_0=1879.65-0.01=1879.64\,\mathrm{mm}
\end{equation*}

\problem{例2-2}
\timu 用例2-1数据，对计算结果进行校核。

\jieda 因 $n$ 为偶数，$n/2=10/2=5$，且 $A=0.01$，由表知
\begin{equation*}
    \sum_{i=1}^{10}v_i=0.01<\frac{n}{2}A=0.05
\end{equation*}
故计算结果正确。

\problem{例2-3}
\timu 测量某量 $11$ 次，求算术平均值并进行校核。

\begin{center}
\scriptsize
\setlength{\tabcolsep}{7pt}
\begin{tabular}{c|ccccccccccc}
\toprule
序号 & 1 & 2 & 3 & 4 & 5 & 6 & 7 & 8 & 9 & 10 & 11\\
\midrule
$l_i/\mathrm{mm}$ & 2000.07 & 2000.09 & 2000.09 & 2000.06 & 2000.07 & 2000.07 & 2000.06 & 2000.08 & 2000.08 & 2000.06 & 2000.07\\
$v_i/\mathrm{mm}$ & $+0.003$ & $+0.023$ & $+0.023$ & $-0.007$ & $+0.003$ & $+0.003$ & $-0.007$ & $+0.013$ & $+0.013$ & $-0.007$ & $+0.003$\\
\bottomrule
\end{tabular}
\end{center}

\jieda 算术平均值为
\begin{equation*}
    \bar{x}=\frac{\sum_{i=1}^{11}l_i}{11}=\frac{22000.74}{11}=2000.0673\,\mathrm{mm}\approx2000.067\,\mathrm{mm}
\end{equation*}
用第一种规则校核：
\begin{equation*}
    \sum_{i=1}^{11}l_i=22000.74\,\mathrm{mm}>11\times2000.067\,\mathrm{mm}=22000.737\,\mathrm{mm}
\end{equation*}
用第二种规则校核：
\begin{equation*}
    \sum_{i=1}^{11}v_i=0.003\,\mathrm{mm}<\left(\frac{11}{2}-0.5\right)A=0.005\,\mathrm{mm}
\end{equation*}
故用两种规则校核均说明计算结果正确。

\problem{例2-4}
\timu 用游标卡尺对某一尺寸测量 $10$ 次，假定已消除系统误差和粗大误差，得到数据如下表。求算术平均值及其标准差。

\begin{center}
\scriptsize
\setlength{\tabcolsep}{6pt}
\begin{tabular}{c|cccccccccc}
\toprule
序号 & 1 & 2 & 3 & 4 & 5 & 6 & 7 & 8 & 9 & 10\\
\midrule
$l_i/\mathrm{mm}$ & 75.01 & 75.04 & 75.07 & 75.00 & 75.03 & 75.09 & 75.06 & 75.02 & 75.05 & 75.08\\
$v_i/\mathrm{mm}$ & $-0.035$ & $-0.005$ & $+0.025$ & $-0.045$ & $-0.015$ & $+0.045$ & $+0.015$ & $-0.025$ & $+0.005$ & $+0.035$\\
$v_i^2/\mathrm{mm}^2$ & 0.001225 & 0.000025 & 0.000625 & 0.002025 & 0.000225 & 0.002025 & 0.000225 & 0.000625 & 0.000025 & 0.001225\\
\bottomrule
\end{tabular}
\end{center}

\jieda 由表得 $\bar{x}=75.045\,\mathrm{mm}$，$\sum v_i=0$，$\sum v_i^2=0.00825\,\mathrm{mm}^2$。根据贝塞尔公式及算术平均值标准差公式，有
\begin{equation*}
    \sigma=\sqrt{\frac{\sum_{i=1}^{n} v_i^2}{n-1}}=\sqrt{\frac{0.00825}{10-1}}=0.0303\,\mathrm{mm},\qquad
    \sigma_{\bar{x}}=\frac{\sigma}{\sqrt{n}}=\frac{0.0303}{\sqrt{10}}=0.0096\,\mathrm{mm}
\end{equation*}
或然误差和平均误差分别为
\begin{equation*}
    R=0.6745\sigma_{\bar{x}}=0.0065\,\mathrm{mm},\qquad
    T=0.7979\sigma_{\bar{x}}=0.0076\,\mathrm{mm}
\end{equation*}

\problem{例2-5}
\timu 仍用例2-4的测量数据，用别捷尔斯法求得标准差。

\jieda 由例2-4数据得 $\sum |v_i|=0.250\,\mathrm{mm}$，故
\begin{equation*}
    \sigma=1.253\frac{\sum_{i=1}^{10}|v_i|}{\sqrt{10(10-1)}}
    =1.253\frac{0.250}{\sqrt{10\times9}}=0.0330\,\mathrm{mm}
\end{equation*}
算术平均值的标准差为
\begin{equation*}
    \sigma_{\bar{x}}=\frac{0.0330}{\sqrt{10}}=0.0104\,\mathrm{mm}
\end{equation*}

\problem{例2-6}
\timu 仍用例2-4的测量数据，用极差法求得标准差。

\jieda 极差为 $\omega_n=75.09-75.00=0.09\,\mathrm{mm}$。当 $n=10$ 时，查表得 $d_n=3.08$，故
\begin{equation*}
    \sigma=\frac{\omega_n}{d_n}=\frac{0.09}{3.08}=0.0292\,\mathrm{mm}
\end{equation*}

\problem{例2-7}
\timu 仍用例2-4的测量数据，用最大误差法求得标准差。

\jieda 由例2-4数据得 $|v_i|_{\max}=0.045\,\mathrm{mm}$。由表查得 $1/K_n=0.57$，则
\begin{equation*}
    \sigma=\frac{|v_i|_{\max}}{K_n}=0.57\times0.045=0.0256\,\mathrm{mm}
\end{equation*}

\problem{例2-8}
\timu 某激光管发出的氖光波长检定为 $\lambda=0.63299130\,\mu\mathrm{m}$。若用来复现更准确的方法测得该波长为 $\lambda=0.63299144\,\mu\mathrm{m}$，试求测定波长的标准差。

\jieda 因原测得的波长是平均值，故可认为其测得值为实际波长，并采用最大误差法估计标准差。取 $1/K_i=1.25$，则系统误差为
\begin{equation*}
    \delta=0.63299130-0.63299144=-14\times10^{-8}\,\mu\mathrm{m}
\end{equation*}
故标准差为
\begin{equation*}
    \sigma=\frac{|\delta|}{K_i}=1.25\times14\times10^{-8}=1.75\times10^{-7}\,\mu\mathrm{m}
\end{equation*}

\newpage

\problem{例2-9}
\timu 对某量进行了 $6$ 次测量，测得数据为 $802.40,802.50,802.38,802.48,802.42,802.46$，求算术平均值及其极限误差。

\jieda 算术平均值、标准差及算术平均值标准差为
\begin{equation*}
    \bar{x}=802.44,
    \qquad
    \sigma=0.047,
    \qquad
    \sigma_{\bar{x}}=\frac{0.047}{\sqrt{6}}=0.019
\end{equation*}
因测量次数较少，应按 $t$ 分布计算算术平均值的极限误差。已知 $n=6$，$\nu=n-1=5$，取 $\alpha=0.01$，查表得 $t=4.03$，故
\begin{equation*}
    \delta_{\lim}\bar{x}=\pm t\sigma_{\bar{x}}=\pm4.03\times0.019=\pm0.076
\end{equation*}

\problem{例2-10}
\timu 对一线值长度的长短进行了三组不等精度测量，其结果为
$x_1=2000.45\,\mathrm{mm}$，$\sigma_1=0.05\,\mathrm{mm}$；
$x_2=2000.15\,\mathrm{mm}$，$\sigma_2=0.20\,\mathrm{mm}$；
$x_3=2000.60\,\mathrm{mm}$，$\sigma_3=0.10\,\mathrm{mm}$。求各测量结果的权。

\jieda 由权与方差成反比可得
\begin{equation*}
    p_1:p_2:p_3=\frac{1}{0.05^2}:\frac{1}{0.20^2}:\frac{1}{0.10^2}=16:1:4
\end{equation*}
因比值的权可取为 $p_1=16$，$p_2=1$，$p_3=4$。

\problem{例2-11}
\timu 三个工作基准米尺在实验室中与国家基准器相比较，得到工作基准米尺的平均长度分别为 $999.9425\,\mathrm{mm}$、$999.9416\,\mathrm{mm}$、$999.9419\,\mathrm{mm}$，其权分别为 $3,2,5$，求加权算术平均值。

\jieda 选取参考值 $999.9400\,\mathrm{mm}$，有
\begin{equation*}
    \bar{x}=999.9400\,\mathrm{mm}+\frac{3\times0.0025+2\times0.0016+5\times0.0019}{3+2+5}
    =999.9420\,\mathrm{mm}
\end{equation*}

\problem{例2-12}
\timu 用 A、B 两种电压表对 $5\,\mathrm{V}$ 电压进行测量。A 表测得 $U_A=5.005\,\mathrm{V}$，$\sigma_A=0.006\,\mathrm{V}$；B 表测得 $U_B=5.002\,\mathrm{V}$，$\sigma_B=0.008\,\mathrm{V}$。求该电压的最佳估计值。

\jieda 按照测量结果标准差确定权，有
\begin{equation*}
    p_A:p_B=\frac{1}{\sigma_A^2}:\frac{1}{\sigma_B^2}
    =\frac{1}{0.006^2}:\frac{1}{0.008^2}=16:9
\end{equation*}
故最佳估计值为
\begin{equation*}
    \bar{U}=\frac{5.005\times16+5.002\times9}{16+9}=5.004\,\mathrm{V}
\end{equation*}

\problem{例2-13}
\timu 求例2-11的加权算术平均值的标准差。

\jieda 由加权算术平均值 $\bar{x}=999.9420\,\mathrm{mm}$，可得各组测量结果的残余误差为
\begin{equation*}
    v_{\bar{x}_1}=+0.5\,\mu\mathrm{m},\qquad
    v_{\bar{x}_2}=-0.4\,\mu\mathrm{m},\qquad
    v_{\bar{x}_3}=-0.1\,\mu\mathrm{m}
\end{equation*}
已知 $m=3$，$p_1=3$，$p_2=2$，$p_3=5$，代入式得
\begin{equation*}
    \sigma_{\bar{x}}
    =\sqrt{\frac{3\times0.5^2+2\times(-0.4)^2+5\times(-0.1)^2}{(3-1)(3+2+5)}}\,\mu\mathrm{m}
    =\sqrt{\frac{1.12}{20}}\,\mu\mathrm{m}=0.24\,\mu\mathrm{m}\approx0.0002\,\mathrm{mm}
\end{equation*}

\problem{例2-14}
\timu 求例2-12的加权算术平均值的标准差。

\jieda 按照例2-12中两个测量值选取的权 $p_A=16$ 和 $p_B=9$，以及测量结果 $5.005\,\mathrm{V}$ 的标准差 $\sigma_A=0.006\,\mathrm{V}$，可得
\begin{equation*}
    \sigma_{\bar{V}}=\sigma_A\sqrt{\frac{p_A}{\sum_{i=1}^{2}p_i}}
    =0.006\sqrt{\frac{16}{16+9}}\,\mathrm{V}=0.0048\,\mathrm{V}\approx0.005\,\mathrm{V}
\end{equation*}

\problem{例2-15}
\timu 对恒温箱温度测量 $10$ 次，测得数据见下表。

\begin{center}
\scriptsize
\setlength{\tabcolsep}{8pt}
\begin{tabular}{c|cccccccccc}
\toprule
序号 & 1 & 2 & 3 & 4 & 5 & 6 & 7 & 8 & 9 & 10\\
\midrule
$t_i/{}^\circ\mathrm{C}$ & 20.06 & 20.07 & 20.06 & 20.08 & 20.10 & 20.12 & 20.14 & 20.18 & 20.18 & 20.21\\
$v_i/{}^\circ\mathrm{C}$ & $-0.06$ & $-0.05$ & $-0.06$ & $-0.04$ & $-0.02$ & 0 & $+0.02$ & $+0.06$ & $+0.06$ & $+0.09$\\
\bottomrule
\end{tabular}
\end{center}

\jieda 由表得 $\bar{x}=20.12{}^\circ\mathrm{C}$，$\sigma=0.055{}^\circ\mathrm{C}$，且
\begin{equation*}
    \sum_{i=1}^{5}v_i=-0.23{}^\circ\mathrm{C},\qquad
    \sum_{i=6}^{10}v_i=+0.23{}^\circ\mathrm{C},\qquad
    \sum_{i=1}^{10}v_i=0
\end{equation*}
由表和残余误差图可知，残余误差符号由负变正，误差值由小到大，故测量列中存在线性系统误差。

\problem{例2-16}
\timu 仍用例2-15的测量数据，按残余误差校核法判断是否含有线性系统误差。

\jieda 此时 $n=10$，$K=5$，而
\begin{equation*}
    \Delta=(-0.06-0.05-0.06-0.04-0.02){}^\circ\mathrm{C}
    -(0+0.02+0.06+0.06+0.09){}^\circ\mathrm{C}
    =-0.46{}^\circ\mathrm{C}
\end{equation*}
因为差值 $\Delta$ 显著不为零，故测量列中含有线性系统误差。

\problem{例2-17}
\timu 雷莱用不同方法制取氮，测得氮气相对密度平均值及其标准差为：由化学法制取氮，$\bar{x}_1=2.29971$，$\sigma_1=0.00041$；由大气中提取氮，$\bar{x}_2=2.31022$，$\sigma_2=0.00019$。判断两种方法间是否存在系统误差。

\jieda 两者差值为
\begin{equation*}
    \Delta=\bar{x}_2-\bar{x}_1=0.01051
\end{equation*}
而标准差为
\begin{equation*}
    \sigma=\sqrt{\sigma_1^2+\sigma_2^2}
    =\sqrt{0.00041^2+0.00019^2}=0.00045
\end{equation*}
因
\begin{equation*}
    \Delta\gg2\sqrt{\sigma_1^2+\sigma_2^2}=2\times0.00045=0.0009
\end{equation*}
故两种方法间存在系统误差。

\problem{例2-18}
\timu 对某量测得两组数据如下，判断两组间有无系统误差。
\begin{equation*}
    x_i:14.7,	14.8,	15.2,	15.6,
    \qquad
    y_i:14.6,	15.0,	15.1
\end{equation*}

\jieda 将两组数据混合后按大小顺序排列，得到秩次：$y_1$ 为 $1$，$x_1$ 为 $2$，$x_2$ 为 $3$，$y_2$ 为 $4$，$y_3$ 为 $5$，$x_3$ 为 $6$，$x_4$ 为 $7$。已知 $n_1=3$，$n_2=4$，计算秩和
\begin{equation*}
    T=1+4+5=10
\end{equation*}
查表得 $T_-=7$，$T_+=17$，因
\begin{equation*}
    T_-=7<T=10<17=T_+
\end{equation*}
故无根据怀疑两组间存在系统误差。

\newpage

\problem{例2-19}
\timu 对某量测得两组数据
\begin{equation*}
\begin{gathered}
    x:1.9,	0.8,	1.1,	0.1,	-0.1,	4.4,	5.5,	1.6,	4.6,	3.4,\\
    y:0.7,	-1.6,	-0.2,	-1.2,	-0.1,	3.4,	3.7,	0.8,	0.0,	2.0
\end{gathered}
\end{equation*}
用 $t$ 检验法判断两组间有无系统误差。

\jieda 计算得
\begin{equation*}
    \bar{x}=2.33,	\bar{y}=0.75,	\sigma_x^2=3.61,	\sigma_y^2=2.89
\end{equation*}
则
\begin{equation*}
    t=(2.33-0.75)\sqrt{\frac{10\times10\times(10+10-2)}{(10+10)(10\times3.61+10\times2.89)}}=1.86
\end{equation*}
由 $\nu=10+10-2=18$ 及 $\alpha=0.05$，查 $t$ 分布表得 $t_a=2.10$。因 $|t|=1.86<t_a=2.10$，故无根据怀疑两组间有系统误差。

\problem{例2-20}
\timu 对某量进行 $15$ 次等精度测量，测得值如下表所列。设这些测得值已消除了系统误差，试判断该测量列中是否含有粗大误差的测得值。

\begin{center}
\scriptsize
\setlength{\tabcolsep}{7pt}
\begin{tabular}{c|ccccc}
\toprule
序号 & 1 & 2 & 3 & 4 & 5\\
\midrule
$l$ & 20.42 & 20.43 & 20.40 & 20.43 & 20.42\\
$v$ & $+0.016$ & $+0.026$ & $-0.004$ & $+0.026$ & $+0.016$\\
\midrule
序号 & 6 & 7 & 8 & 9 & 10\\
\midrule
$l$ & 20.43 & 20.39 & 20.30 & 20.40 & 20.43\\
$v$ & $+0.026$ & $-0.014$ & $-0.104$ & $-0.004$ & $+0.026$\\
\midrule
序号 & 11 & 12 & 13 & 14 & 15\\
\midrule
$l$ & 20.42 & 20.41 & 20.39 & 20.39 & 20.40\\
$v$ & $+0.016$ & $+0.006$ & $-0.014$ & $-0.014$ & $-0.004$\\
\bottomrule
\end{tabular}
\end{center}

\jieda 由表得
\begin{equation*}
    \bar{x}=20.404,	\sigma=\sqrt{\frac{0.0150}{14}}=0.033,	3\sigma=0.099
\end{equation*}
根据 $3\sigma$ 准则，第八个测得值的残余误差为
\begin{equation*}
    |v_8|=0.104>0.099
\end{equation*}
故第八个测得值含有粗大误差，应予剔除。删除后重新计算得
\begin{equation*}
    \bar{x}'=20.411,	\sigma'=\sqrt{\frac{0.00337}{13}}=0.016,	3\sigma'=0.048
\end{equation*}
剩下的 $14$ 个测得值的残余误差均满足 $|v_i'|<3\sigma'$，故可认为这些测得值不再含有粗大误差。

\problem{例2-21}
\timu 试判断例2-20中是否含有粗大误差。

\jieda 首先怀疑第八测得值含有粗大误差，将其剔除。然后根据表2-11中剩下的 $14$ 个测量值计算平均值和标准差，得
\begin{equation*}
    \bar{x}=20.411,	\sigma=0.016
\end{equation*}
选取显著度 $\alpha=0.05$，已知 $n=15$，查表得 $K(15,0.05)=2.24$，则
\begin{equation*}
    K\sigma=2.24\times0.016=0.036
\end{equation*}
因
\begin{equation*}
    |x_8-\bar{x}|=|20.30-20.411|=0.111>0.036
\end{equation*}
故第八测量值含有粗大误差，应予剔除。然后对剩下 $14$ 个测量值进行判断，可知这些测量值不再含有粗大误差。

\problem{例2-22}
\timu 用例2-20测得值，试判断该测量列中的测得值是否含有粗大误差。

\jieda 计算得 $\bar{x}=20.404$，$\sigma=0.033$。按测得值大小顺序排列有 $x_{(1)}=20.30$，$x_{(15)}=20.43$。首先判断 $x_{(1)}$ 是否含有粗大误差：
\begin{equation*}
    g_{(1)}=\frac{20.404-20.30}{0.033}=3.15
\end{equation*}
查表得 $g_0(15,0.05)=2.41$，故
\begin{equation*}
    g_{(1)}=3.15>g_0(15,0.05)=2.41
\end{equation*}
所以第八个测得值含有粗大误差，应予剔除。对剩下 $14$ 个数据再判别，得 $g_{(15)}=1.18<g_0(14,0.05)=2.37$，故其余测得值不含粗大误差。

\newpage

\problem{例2-23}
\timu 同例2-20测量数据，将 $x_i$ 排成顺序量，按狄克松准则判断粗大误差。

\begin{center}
\scriptsize
\setlength{\tabcolsep}{5pt}
\begin{tabular}{c|ccccccccccccccc}
\toprule
$x_i$ & 20.30 & 20.39 & 20.39 & 20.39 & 20.40 & 20.40 & 20.40 & 20.41 & 20.42 & 20.42 & 20.42 & 20.43 & 20.43 & 20.43 & 20.43\\
顺序号 $x_{(i)}$ & 1 & 2 & 3 & 4 & 5 & 6 & 7 & 8 & 9 & 10 & 11 & 12 & 13 & 14 & 15\\
顺序号 $x'_{(i)}$ & -- & 1 & 2 & 3 & 4 & 5 & 6 & 7 & 8 & 9 & 10 & 11 & 12 & 13 & 14\\
\bottomrule
\end{tabular}
\end{center}

\jieda 首先判断最大值 $x_{(15)}$。因 $n=15$，按式计算统计量
\begin{equation*}
    r_{22}=\frac{x_{(15)}-x_{(13)}}{x_{(15)}-x_{(3)}}=\frac{20.43-20.43}{20.43-20.39}=0
\end{equation*}
查表得 $r_0(15,0.05)=0.525$，故 $x_{(15)}$ 不含粗大误差。再判断最小值 $x_{(1)}$：
\begin{equation*}
    r_{22}=\frac{x_{(1)}-x_{(3)}}{x_{(1)}-x_{(13)}}=\frac{20.30-20.39}{20.30-20.43}=0.692
\end{equation*}
因 $r_{22}>r_0=0.525$，故 $x_{(1)}$ 含有粗大误差，应予剔除。剩下 $14$ 个数据再判别，可认为不再含有粗大误差。

\problem{例2-24}

\timu 对某一轴径等精度测量 $9$ 次，得到下表数据，求测量结果。

\begin{center}
\scriptsize
\setlength{\tabcolsep}{6pt}
\begin{tabular}{c|ccccccccc}
\toprule
序号 & 1 & 2 & 3 & 4 & 5 & 6 & 7 & 8 & 9\\
\midrule
$l_i/\mathrm{mm}$ & 24.774 & 24.778 & 24.771 & 24.780 & 24.772 & 24.777 & 24.773 & 24.775 & 24.774\\
$v_i/\mathrm{mm}$ & $-0.001$ & $+0.003$ & $-0.004$ & $+0.005$ & $-0.003$ & $+0.002$ & $-0.002$ & 0 & $-0.001$\\
$v_i^2/\mathrm{mm}^2$ & 0.000001 & 0.000009 & 0.000016 & 0.000025 & 0.000009 & 0.000004 & 0.000004 & 0 & 0.000001\\
\bottomrule
\end{tabular}
\end{center}

\jieda 假定该测量列不存在固定的系统误差，则可按下列步骤求测量结果。

\begin{enumerate}[label=\arabic*.]
    \item 求算术平均值

    测量列的算术平均值 $\bar{x}$ 为
    \begin{equation*}
        \bar{x}
        =
        \frac{\displaystyle\sum_{i=1}^{9}l_i}{n}
        =
        \frac{222.974}{9}\,\mathrm{mm}
        =
        24.7749\,\mathrm{mm}
        \approx24.775\,\mathrm{mm}
    \end{equation*}

    \item 求残余误差

    各测得值的残余误差 $v_i=l_i-\bar{x}$，并列入表中。
\newpage
    \item 校核算术平均值及其残余误差

    根据残余误差代数和校核规则，现用规则 $2$ 进行校核，因
    \begin{equation*}
        A=0.001\,\mathrm{mm},\qquad n=9
    \end{equation*}
    由上表知
    \begin{equation*}
        \left|\sum_{i=1}^{9}v_i\right|
        =
        0.001\,\mathrm{mm}
        <
        \left(\frac{n}{2}-0.5\right)A
        =
        4\times0.001\,\mathrm{mm}
        =
        0.004\,\mathrm{mm}
    \end{equation*}
    故以上计算正确。若发现计算有误，应重新进行上述计算和校核。

    \item 判断系统误差

    根据残余误差观察法，由上表可以看出误差符号大体上正负相间，且无显著变化规律，因此可判断该测量列无变化的系统误差存在。

    若按残余误差校核法，因 $n=9$，则
    \begin{equation*}
        K=\frac{n+1}{2}=5
    \end{equation*}
    \begin{equation*}
        \Delta
        =
        \sum_{i=1}^{5}v_i-\sum_{j=6}^{9}v_j
        =
        [0-(-0.001)]\,\mathrm{mm}
        =
        0.001\,\mathrm{mm}
    \end{equation*}
    因差值 $A$ 较小，故也可判断该测量列无变化的系统误差存在。

    \item 求测量列单次测量的标准差

    根据贝塞尔公式或别捷尔斯公式，求得测量列单次测量的标准差 $\sigma$ 为
    \begin{equation*}
        \sigma
        =
        \sqrt{\frac{\displaystyle\sum_{i=1}^{9}v_i^2}{n-1}}
        =
        \sqrt{\frac{0.000069}{8}}\,\mathrm{mm}
        =
        0.0029\,\mathrm{mm}
    \end{equation*}
    \begin{equation*}
        \sigma'
        =
        1.253
        \frac{\displaystyle\sum_{i=1}^{n}|v_i|}
        {\sqrt{n(n-1)}}
        =
        1.253\times
        \frac{0.021}{\sqrt{9\times8}}\,\mathrm{mm}
        =
        0.0031\,\mathrm{mm}
    \end{equation*}
    用两种方法计算的标准差比值为
    \begin{equation*}
        \frac{\sigma'}{\sigma}
        =
        \frac{0.0031}{0.0029}
        =
        1.069
        =
        1+u
    \end{equation*}
    \begin{equation*}
        u=0.069
    \end{equation*}
    因
    \begin{equation*}
        |u|
        =
        0.069
        <
        \frac{2}{\sqrt{n-1}}
        =
        \frac{2}{\sqrt{8}}
        \approx0.707
    \end{equation*}
    故同样可判断该测量列无系统误差存在。

    \item 判别粗大误差

    根据 $3\sigma$ 判别准则的适用特点，本实例测量轴径的次数较少，因而不采用 $3\sigma$ 准则来判别粗大误差。

    若按格罗布斯判别准则，将测得值按大小顺序排列后有
    \begin{equation*}
        x_{(1)}=24.771\,\mathrm{mm},
        \qquad
        x_{(9)}=24.780\,\mathrm{mm}
    \end{equation*}
    \begin{equation*}
        x_{(2)}-x_{(1)}
        =
        24.775\,\mathrm{mm}
        -
        24.771\,\mathrm{mm}
        =
        0.004\,\mathrm{mm}
    \end{equation*}
    \begin{equation*}
        x_{(9)}-x_{(8)}
        =
        24.780\,\mathrm{mm}
        -
        24.775\,\mathrm{mm}
        =
        0.005\,\mathrm{mm}
    \end{equation*}

    首先判断 $x_{(9)}$ 是否含有粗大误差：
    \begin{equation*}
        g_{(9)}
        =
        \frac{24.780-24.775}{0.0029}
        =
        1.72
    \end{equation*}
    查表得
    \begin{equation*}
        g_{(9)}(9,0.05)=2.11
    \end{equation*}
    因
    \begin{equation*}
        g_{(9)}=1.70<g_{(9)}(9,0.05)=2.11,
        \qquad
        \text{且}\quad
        g_{(1)}<g_{(9)}
    \end{equation*}
    故可判断测量列不存在粗大误差。

    若发现测量列存在粗大误差，应将含有粗大误差的测得值剔除，然后再按上述步骤重新计算，直至所有测得值皆不包含粗大误差时为止。

    \item 求算术平均值的标准差

    计算 $\sigma_{\bar{x}}$ 得
    \begin{equation*}
        \sigma_{\bar{x}}
        =
        \frac{\sigma}{\sqrt{n}}
        =
        \frac{0.0029}{\sqrt{9}}\,\mathrm{mm}
        \approx0.001\,\mathrm{mm}
    \end{equation*}

    \item 求算术平均值的极限误差

    因为测量列的测量次数较少，算术平均值的极限误差按 $t$ 分布计算。

    已知 $\nu=n-1=8$，取 $\alpha=0.05$，查表得
    \begin{equation*}
        t_{\alpha}=2.31
    \end{equation*}
    求得算术平均值的极限误差 $\delta_{\lim}\bar{x}$ 为
    \begin{equation*}
        \delta_{\lim}\bar{x}
        =
        \pm t_{\alpha}\sigma_{\bar{x}}
        =
        \pm2.31\times0.001\,\mathrm{mm}
        =
        \pm0.0023\,\mathrm{mm}
    \end{equation*}

    \item 写出最后测量结果

    最后测量结果通常用算术平均值及其极限误差来表示，即
    \begin{equation*}
        L
        =
        \bar{x}+\delta_{\lim}\bar{x}
        =
        (24.775\pm0.0023)\,\mathrm{mm}
    \end{equation*}
\end{enumerate}

\newpage

\problem{例2-25}
\timu 对某一角度进行 $6$ 组不等精度测量，各组的单次测量均为等精度测量，其测量结果为：
\begin{equation*}
\begin{aligned}
    \text{测 }6\text{ 次得 }\alpha_1
    &=75^\circ18'06'',\qquad
    \text{测 }30\text{ 次得 }\alpha_2
    &=75^\circ18'10'',\\
    \text{测 }24\text{ 次得 }\alpha_3
    &=75^\circ18'08'',\qquad
    \text{测 }12\text{ 次得 }\alpha_4
    &=75^\circ18'16'',\\
    \text{测 }12\text{ 次得 }\alpha_5
    &=75^\circ18'13'',\qquad
    \text{测 }36\text{ 次得 }\alpha_6
    &=75^\circ18'09''
\end{aligned}
\end{equation*}
求最后测量结果。

\jieda 根据测量次数确定各组权：
\begin{equation*}
    p_1:p_2:p_3:p_4:p_5:p_6=1:5:4:2:2:6,
    \qquad
    \sum_{i=1}^{6}p_i=20
\end{equation*}
选取参考值 $\alpha_0=75^\circ18'06''$，则
\begin{equation*}
    \bar{\alpha}=75^\circ18'06''+\frac{1\times0''+5\times4''+4\times2''+2\times10''+2\times7''+6\times3''}{20}=75^\circ18'10''
\end{equation*}
残余误差为
\begin{equation*}
    v_1=-4'',	 v_2=0,	 v_3=-2'',	 v_4=6'',	 v_5=3'',	 v_6=-1''
\end{equation*}
并有 $\sum p_iv_i=0$，故计算正确。加权算术平均值标准差为
\begin{equation*}
    \sigma_{\bar{x}}=
    \sqrt{\frac{1\times(4'')^2+5\times0''+4\times(2'')^2+2\times(6'')^2+2\times(3'')^2+6\times(1'')^2}{(6-1)\times20}}
    =1.1''
\end{equation*}
取置信系数 $t=3$，则极限误差为
\begin{equation*}
    \delta_{\lim}\bar{x}=\pm3\sigma_{\bar{x}}=\pm3.3''
\end{equation*}
最后测量结果为
\begin{equation*}
    \alpha=\bar{\alpha}+\delta_{\lim}\bar{x}=75^\circ18'10''\pm3.3''
\end{equation*}

\newpage

\problem{例2-26}
\timu 电子电量 $e$ 与质量 $m$ 比值 $e/m$ 的两次观测结果为
\begin{equation*}
    x_1\pm\sigma_1=1.75080\pm0.00042,
    \qquad
    x_2\pm\sigma_2=1.75059\pm0.00036
\end{equation*}
单位为 $10^{11}\,\mathrm{C/kg}$。两次观测结果不存在系统误差和粗大误差，求最后的测量结果。

\jieda 首先根据两次测量的标准差求权：
\begin{equation*}
    p_1:p_2=\frac{1}{\sigma_1^2}:\frac{1}{\sigma_2^2}
    =\frac{1}{0.00042^2}:\frac{1}{0.00036^2}=36:49
\end{equation*}
取 $p_1=36$，$p_2=49$，则加权算术平均值为
\begin{equation*}
    \bar{x}=\frac{x_1p_1+x_2p_2}{p_1+p_2}
    =\frac{1.75080\times36+1.75059\times49}{36+49}\times10^{11}\,\mathrm{C/kg}
    =1.75068\times10^{11}\,\mathrm{C/kg}
\end{equation*}
残余误差为
\begin{equation*}
    v_1=0.00012\times10^{11}\,\mathrm{C/kg},
    \qquad
    v_2=-0.00009\times10^{11}\,\mathrm{C/kg}
\end{equation*}
加权平均值的标准差为
\begin{equation*}
    \sigma_{\bar{x}}=\sigma_{x_1}\sqrt{\frac{p_1}{p_1+p_2}}
    =0.00042\sqrt{\frac{36}{36+49}}\times10^{11}\,\mathrm{C/kg}
    =0.00027\times10^{11}\,\mathrm{C/kg}
\end{equation*}
若取置信系数 $t=3$，则最后极限误差为
\begin{equation*}
    \delta_{\lim}\bar{x}=\pm3\sigma_{\bar{x}}=\pm0.00081\times10^{11}\,\mathrm{C/kg}
\end{equation*}
最后测量结果为
\begin{equation*}
    x=\bar{x}+\delta_{\lim}\bar{x}=(1.75068\times10^{11}\pm0.00081\times10^{11})\,\mathrm{C/kg}
\end{equation*}

\section*{第三章\quad 误差的合成与分配}

\problem{例3-1}
\timu 用弓高弦长法间接测量大直径 $D$，直接测得弓高 $h$ 和弦长 $s$。若 $h=50\,\mathrm{mm}$，$\Delta h=-0.1\,\mathrm{mm}$，$s=500\,\mathrm{mm}$，$\Delta s=1\,\mathrm{mm}$，求测量结果。

\jieda 由几何关系得
\begin{equation*}
    D=\frac{s^2}{4h}+h
\end{equation*}
若不考虑测得值的系统误差，则
\begin{equation*}
    D_0=\frac{500^2}{4\times50}\,\mathrm{mm}+50\,\mathrm{mm}=1300\,\mathrm{mm}
\end{equation*}
直径 $D$ 的系统误差为
\begin{equation*}
    \Delta D=\frac{\partial f}{\partial s}\Delta s+\frac{\partial f}{\partial h}\Delta h
    =\frac{s}{2h}\Delta s-\left(\frac{s^2}{4h^2}-1\right)\Delta h
\end{equation*}
其中
\begin{equation*}
    \frac{\partial f}{\partial s}=\frac{500}{2\times50}=5,
    \qquad
    \frac{\partial f}{\partial h}=-\left(\frac{500^2}{4\times50^2}-1\right)=-24
\end{equation*}
故
\begin{equation*}
    \Delta D=5\times1\,\mathrm{mm}-24\times(-0.1)\,\mathrm{mm}=7.4\,\mathrm{mm}
\end{equation*}
通过修正可消除直径的系统误差，故被测直径的实际尺寸为
\begin{equation*}
    D=D_0-\Delta D=1300\,\mathrm{mm}-7.4\,\mathrm{mm}=1292.6\,\mathrm{mm}
\end{equation*}

\problem{例3-2}
\timu 用某直流电桥测量电阻 $R_x$。当三个桥臂平衡时，$R_1=100.0\,\Omega$，$R_2=50.0\,\Omega$，$R_3=25.0\,\Omega$，且 $\Delta R_1=0.2\,\Omega$，$\Delta R_2=0.1\,\Omega$，$\Delta R_3=0.2\,\Omega$，求电阻 $R_x$ 的测量结果。

\jieda 电桥平衡时
\begin{equation*}
    R_x=\frac{R_1}{R_2}R_3
\end{equation*}
不考虑测得值的系统误差时，
\begin{equation*}
    R_0=\frac{100}{50}\times25\,\Omega=50\,\Omega
\end{equation*}
误差传递系数为
\begin{equation*}
    \frac{\partial f}{\partial R_1}=\frac{R_3}{R_2}=0.5,
    \qquad
    \frac{\partial f}{\partial R_2}=-\frac{R_1R_3}{R_2^2}=-1,
    \qquad
    \frac{\partial f}{\partial R_3}=\frac{R_1}{R_2}=2
\end{equation*}
所以
\begin{equation*}
    \Delta R_x=0.5\times0.2-1\times0.1+2\times0.2=0.4\,\Omega
\end{equation*}
将系统误差修正后，被测电阻为
\begin{equation*}
    R_x=R_0-\Delta R_x=50\,\Omega-0.4\,\Omega=49.6\,\Omega
\end{equation*}

\problem{例3-3}
\timu 对例3-1用弓高弦长法测量大直径 $D$，已知
$h=50\,\mathrm{mm}$，$s=500\,\mathrm{mm}$，$\delta_{\lim}s=\pm0.1\,\mathrm{mm}$，$\delta_{\lim}h=\pm0.05\,\mathrm{mm}$，求直径的极限误差。

\jieda 由例3-1知 $D=s^2/(4h)+h$，根据极限误差合成公式，有
\begin{equation*}
\begin{aligned}
    \delta_{\lim}D
    &=\pm\sqrt{\left(\frac{\partial f}{\partial s}\right)^2\delta_{\lim}^2s+
    \left(\frac{\partial f}{\partial h}\right)^2\delta_{\lim}^2h}\\
    &=\pm\sqrt{\left(\frac{500}{2\times50}\right)^2\times0.1^2+
    \left(\frac{500^2}{4\times50^2}-1\right)^2\times0.05^2}\,\mathrm{mm}\\
    &=\pm\sqrt{1.69}\,\mathrm{mm}=\pm1.3\,\mathrm{mm}
\end{aligned}
\end{equation*}
故直径的最后结果为
\begin{equation*}
    D=1292.6\,\mathrm{mm}\pm1.3\,\mathrm{mm}
\end{equation*}

\problem{例3-4}
\timu 在例3-2中，若各个电阻测量存在随机误差且互相独立，已知
$\sigma_{R_1}=0.4\,\Omega$，$\sigma_{R_2}=0.2\,\Omega$，$\sigma_{R_3}=0.4\,\Omega$，求 $R_x$ 的标准差。

\jieda 根据函数随机误差合成公式，
\begin{equation*}
\begin{aligned}
    \sigma_{R_x}
    &=\sqrt{\left(\frac{\partial f}{\partial R_1}\right)^2\sigma_{R_1}^2+
    \left(\frac{\partial f}{\partial R_2}\right)^2\sigma_{R_2}^2+
    \left(\frac{\partial f}{\partial R_3}\right)^2\sigma_{R_3}^2}\\
    &=\sqrt{0.5^2\times0.4^2+(-1)^2\times0.2^2+2^2\times0.4^2}\,\Omega\\
    &\approx0.82\,\Omega
\end{aligned}
\end{equation*}
由例3-2结果和 $\sigma_{R_x}$ 可得
\begin{equation*}
    R_x=(50-0.4)\,\Omega\pm3\times0.82\,\Omega\approx49.6\,\Omega\pm2.5\,\Omega
\end{equation*}

\problem{例3-5}
\timu 在万能工具显微镜上用影像法测量某一平面工件的长度共两次，测得 $l_1=50.026\,\mathrm{mm}$，$l_2=50.025\,\mathrm{mm}$。已知工件高度 $H=80\,\mathrm{mm}$，求测量结果及其极限误差。

\begin{center}
\small
\setlength{\tabcolsep}{6pt}
\begin{tabular}{c|c|c|c|c}
\toprule
序号 & 误差因素 & 随机误差$/\mu\mathrm{m}$ & 未定系统误差$/\mu\mathrm{m}$ & 备注\\
\midrule
1 & 阿贝误差 & -- & $\pm1$ & --\\
2 & 光学刻尺刻度误差 & -- & $\pm1.25$ & 加修正值时不计入总误差\\
3 & 温度误差 & -- & $\pm0.35$ & --\\
4 & 读数误差 & $\pm0.8$ & -- & --\\
5 & 瞄准误差 & $\pm1$ & -- & --\\
6 & 光学刻尺检定误差 & -- & $\pm0.5$ & 不加修正值时不计入总误差\\
\bottomrule
\end{tabular}
\end{center}

\jieda 两次测量结果平均值为
\begin{equation*}
    L_0=\frac{1}{2}(l_1+l_2)=\frac{1}{2}(50.026+50.025)\,\mathrm{mm}=50.0255\,\mathrm{mm}
\end{equation*}
查刻度误差修正值 $\Delta=-0.0008\,\mathrm{mm}$，修正后
\begin{equation*}
    L=L_0+\Delta=50.0255\,\mathrm{mm}-0.0008\,\mathrm{mm}=50.0247\,\mathrm{mm}
\end{equation*}
未修正刻尺刻度误差时
\begin{equation*}
    \delta=\pm\sqrt{\frac{1}{2}(1^2+0.8^2)+(1^2+1.25^2+0.35^2)}\,\mu\mathrm{m}
    =\pm1.87\,\mu\mathrm{m}\approx\pm1.9\,\mu\mathrm{m}
\end{equation*}
故 $L_0=(50.0255\pm0.0019)\,\mathrm{mm}$。已修正刻尺刻度误差时
\begin{equation*}
    \delta=\pm\sqrt{\frac{1}{2}(1^2+0.8^2)+(1^2+0.35^2+0.5^2)}\,\mu\mathrm{m}
    =\pm1.48\,\mu\mathrm{m}\approx\pm1.5\,\mu\mathrm{m}
\end{equation*}
故测量结果为
\begin{equation*}
    L=(50.0247\pm0.0015)\,\mathrm{mm}
\end{equation*}

\problem{例3-6}
\timu 用 TC328B 型天平，配用三等标准砝码称一不锈钢球质量，一次称量得钢球质量 $M=14.0040\,\mathrm{g}$，求测量结果的标准差。

\jieda 天平示值变动性所引起的随机误差为 $\sigma_1=0.05\,\mathrm{mg}$。标准砝码误差中 $10\,\mathrm{g}$ 砝码标准差为 $u_{11}=0.4\,\mathrm{mg}$，$2\,\mathrm{g}$ 砝码标准差为 $u_{12}=0.2\,\mathrm{mg}$，则
\begin{equation*}
    u_1=\sqrt{u_{11}^2+2u_{12}^2}=\sqrt{0.4^2+2\times0.2^2}\,\mathrm{mg}\approx0.5\,\mathrm{mg}
\end{equation*}
天平示值误差对应的标准差为
\begin{equation*}
    u_2=2\times0.1\times\frac{40}{100}\times\frac{1}{3}\,\mathrm{mg}\approx0.03\,\mathrm{mg}
\end{equation*}
于是测量结果的总标准差为
\begin{equation*}
    \sigma=\sqrt{\sigma_1^2+u_1^2+u_2^2}
    =\sqrt{0.05^2+0.5^2+0.03^2}\,\mathrm{mg}\approx0.5\,\mathrm{mg}
\end{equation*}
最后测量结果可表示为
\begin{equation*}
    M=14.0040\,\mathrm{g}\pm0.0005\,\mathrm{g}
\end{equation*}

\problem{例3-7}
\timu 测量一圆柱体体积，可间接测量圆柱直径 $D$ 及高度 $h$，根据 $V=\pi D^2h/4$ 求体积 $V$。若要求测量体积的相对误差为 $1\%$，已知 $D_0=20\,\mathrm{mm}$，$h_0=50\,\mathrm{mm}$，试确定 $D$ 和 $h$ 的测量精度。

\jieda 取 $\pi=3.1416$，可得
\begin{equation*}
    V_0=\frac{\pi D_0^2}{4}h_0=\frac{3.1416\times20^2}{4}\times50\,\mathrm{mm}^3=15708\,\mathrm{mm}^3
\end{equation*}
体积绝对误差为 $\delta_V=V_0\times1\%=157.08\,\mathrm{mm}^3$。按等作用原则分配误差，有
\begin{equation*}
    \delta_D=\frac{\delta_V}{\sqrt{2}}\frac{1}{\partial V/\partial D}
    =\frac{157.08}{\sqrt{2}}\frac{2}{\pi\times20\times50}\,\mathrm{mm}=0.071\,\mathrm{mm}
\end{equation*}
\begin{equation*}
    \delta_h=\frac{\delta_V}{\sqrt{2}}\frac{1}{\partial V/\partial h}
    =\frac{157.08}{\sqrt{2}}\frac{4}{\pi\times20^2}\,\mathrm{mm}=0.351\,\mathrm{mm}
\end{equation*}
由此可知，测量直径 $D$ 的精度需要高些，测量高度 $h$ 的精度可以低些。

\newpage

\problem{例3-8}
\timu 测量某箱体零件的轴心距 $L$，有三种方案：$L=L_1-\frac{1}{2}d_1-\frac{1}{2}d_2$，$L=L_2+\frac{1}{2}d_1+\frac{1}{2}d_2$，$L=\frac{1}{2}L_1+\frac{1}{2}L_2$。若 $\sigma_{d1}=5\,\mu\mathrm{m}$，$\sigma_{d2}=7\,\mu\mathrm{m}$，$\sigma_{L1}=8\,\mu\mathrm{m}$，$\sigma_{L2}=10\,\mu\mathrm{m}$，试选择最佳测量方案。

\jieda 三种方法的函数标准差分别为

第一种方法
\begin{equation*}
\begin{aligned}
    \sigma_L
    &=
    \sqrt{
    \left(\frac{\partial f}{\partial L_1}\right)^2\sigma_{L_1}^2
    +
    \left(\frac{\partial f}{\partial d_1}\right)^2\sigma_{d_1}^2
    +
    \left(\frac{\partial f}{\partial d_2}\right)^2\sigma_{d_2}^2
    }\\
    &=
    \sqrt{
    \sigma_{L_1}^2
    +
    \left(\frac{1}{2}\right)^2\sigma_{d_1}^2
    +
    \left(\frac{1}{2}\right)^2\sigma_{d_2}^2
    }\\
    &=
    \sqrt{8^2+2.5^2+3.5^2}\,\mu\mathrm{m}
    =
    9.1\,\mu\mathrm{m}
\end{aligned}
\end{equation*}

第二种方法
\begin{equation*}
\begin{aligned}
    \sigma_L
    &=
    \sqrt{
    \left(\frac{\partial f}{\partial L_2}\right)^2\sigma_{L_2}^2
    +
    \left(\frac{\partial f}{\partial d_1}\right)^2\sigma_{d_1}^2
    +
    \left(\frac{\partial f}{\partial d_2}\right)^2\sigma_{d_2}^2
    }\\
    &=
    \sqrt{
    \sigma_{L_2}^2
    +
    \left(\frac{1}{2}\right)^2\sigma_{d_1}^2
    +
    \left(\frac{1}{2}\right)^2\sigma_{d_2}^2
    }\\
    &=
    \sqrt{10^2+2.5^2+3.5^2}\,\mu\mathrm{m}
    =
    10.9\,\mu\mathrm{m}
\end{aligned}
\end{equation*}

第三种方法
\begin{equation*}
\begin{aligned}
    \sigma_L
    &=
    \sqrt{
    \left(\frac{\partial f}{\partial d_1}\right)^2\sigma_{d_1}^2
    +
    \left(\frac{\partial f}{\partial d_2}\right)^2\sigma_{d_2}^2
    }\\
    &=
    \sqrt{
    \left(\frac{1}{2}\right)^2\sigma_{d_1}^2
    +
    \left(\frac{1}{2}\right)^2\sigma_{d_2}^2
    }\\
    &=
    \sqrt{4^2+5^2}\,\mu\mathrm{m}
    =
    6.4\,\mu\mathrm{m}
\end{aligned}
\end{equation*}
由计算结果可知，第三种方法误差最小，故应选第三种测量方案。

\problem{例3-9}
\timu 对例3-3的弓高弦长法测量直径 $D$，已知 $D=s^2/(4h)+h$，试确定最佳测量方案。

\jieda 直径标准差为
\begin{equation*}
    \sigma_D=\sqrt{\left(\frac{s}{2h}\right)^2\sigma_s^2+
    \left(\frac{s^2}{4h^2}-1\right)^2\sigma_h^2}
\end{equation*}
欲使 $\sigma_D$ 最小，应使误差传递系数尽量小。若使 $s/(2h)=0$，则需 $s=0$，无实际意义；若使 $s^2/(4h^2)-1=0$，则 $s=2h$，即要求测量直径。由此可知，弓高越接近 $s/2$，直径测量误差越小。

\problem{例3-10}
\timu 测量金属导线的电导率 $\gamma$，已知函数式为
\begin{equation*}
    \gamma=\frac{4l}{\pi d^2R}
\end{equation*}
其中 $l,d,R$ 分别为导线长度、直径和电阻，试分析最佳测量方案。

\jieda 电导率标准差为
\begin{equation*}
    \sigma_{\gamma}=\sqrt{
    \left(\frac{4}{\pi d^2R}\right)^2\sigma_l^2+
    \left(\frac{8l}{\pi d^3R}\right)^2\sigma_d^2+
    \left(\frac{4l}{\pi d^2R^2}\right)^2\sigma_R^2}
\end{equation*}
由误差传递系数可知，为减小 $\sigma_{\gamma}$，宜选用长度较小、直径较大的金属导线，即用短而粗的导线测量金属电导率。同时导线直径的标准差对电导率影响较大，故需用高精度方法测量导线直径。

\section*{第四章\quad 测量不确定度}

\problem{例4-1}

\timu 某校准证书说明，标称值 $1\,\mathrm{kg}$ 的标准砝码的质量 $m_s$ 为 $1000.000325\,\mathrm{g}$，该值的测量不确定度按三倍标准差计算为 $240\,\mu\mathrm{g}$，求该砝码质量的标准不确定度。

\jieda 已知测量不确定度 $U_{m_s}=240\,\mu\mathrm{g}$，$k=3$，故标准不确定度为
\begin{equation*}
    u_{m_s}
    =
    \frac{U_{m_s}}{k}
    =
    \frac{240\,\mu\mathrm{g}}{3}
    =
    80\,\mu\mathrm{g}
\end{equation*}

\problem{例4-2}

\timu 由于用查得纯铜在温度 $20^\circ\mathrm{C}$ 时的线膨胀系数 $\alpha$ 为 $16.52\times10^{-6}/^\circ\mathrm{C}$，并已知该系数的误差范围为 $\pm0.4\times10^{-6}/^\circ\mathrm{C}$，求线膨胀系数 $\alpha$ 的标准不确定度。

\jieda 根据手册提供的信息可认为 $\alpha$ 的值以等概率落于区间 $(16.52-0.4)\times10^{-6}/^\circ\mathrm{C}$ 至 $(16.52+0.4)\times10^{-6}/^\circ\mathrm{C}$ 内，且不可能位于此区间之外，故假设 $\alpha$ 服从均匀分布。已知其区间半宽 $a=0.4\times10^{-6}/^\circ\mathrm{C}$，则纯铜在温度为 $20^\circ\mathrm{C}$ 的线膨胀系数 $\alpha$ 的标准不确定度为
\begin{equation*}
    u_{\alpha}
    =
    \frac{a}{\sqrt{3}}
    =
    \frac{0.4\times10^{-6}/^\circ\mathrm{C}}{\sqrt{3}}
    =
    0.23\times10^{-6}/^\circ\mathrm{C}
\end{equation*}

\problem{例4-3}

\timu 用标准数字电压表在标准条件下，对直流电压源 $10\,\mathrm{V}$ 点的输出电压值进行重复测量 $10$ 次，测得值如下：

\begin{center}
\scriptsize
\setlength{\tabcolsep}{6pt}
\begin{tabular}{c|cccccccccc}
\toprule
$n$ & 1 & 2 & 3 & 4 & 5 & 6 & 7 & 8 & 9 & 10\\
\midrule
$V_i$ & 10.000107 & 10.000103 & 10.000097 & 10.000111 & 10.000091 & 10.000108 & 10.000121 & 10.000101 & 10.000110 & 10.000094\\
\bottomrule
\end{tabular}
\end{center}

已知标准电压表的示值误差按 $3$ 倍标准差计算为 $\pm3.5\times10^{-6}\times U_x$（$U_x$ 为标准电压表的读数，相对标准差为 $25\%$），$24\mathrm{h}$ 的稳定度不超过 $\pm15\,\mu\mathrm{V}$（按均匀分布，相对标准差为 $10\%$）。若用 $10$ 次测量的平均值作为电压测量的估计值，试分析电压测量的不确定度来源，并分别计算其标准不确定度和自由度。

\jieda 根据题意，引起电压测量不确定度的主要来源有电压测量的重复性、标准电压表的示值误差及其稳定性，并相应标定不确定度的分量计算如下：

\begin{enumerate}[label=(\arabic*)]
    \item 电压测量重复性引起的标准不确定度 $u_1$

    由 $10$ 次测量数据计算的平均值 $\bar{V}=10.000104\,\mathrm{V}$，用贝塞尔法计算单次测量的标准差为
    \begin{equation*}
        \sigma
        =
        \sqrt{
        \frac{\displaystyle\sum_{i=1}^{10}(V_i-\bar{V})^2}{10-1}
        }
        =
        0.000009\,\mathrm{V}
        =
        9\,\mu\mathrm{V}
    \end{equation*}

    平均值的标准差
    \begin{equation*}
        \sigma_{\bar{V}}
        =
        \frac{\sigma}{\sqrt{10}}
        =
        2.8\,\mu\mathrm{V}
    \end{equation*}

    电压测量重复性引起的标准不确定度 $u_1$ 的计算属 A 类评定，则 $u_1=\sigma_{\bar{V}}=2.8\,\mu\mathrm{V}$，其自由度 $\nu_1=10-1=9$。

    \item 标准电压表示值误差引起的标准不确定度 $u_2$

    标准电压表的示值误差按 $3$ 倍标准差计算，不确定度 $u_2$ 的计算属 B 类评定， $U_x=3.5\times10^{-6}\times10\,\mathrm{V}$，$k=3$，则标准不确定度为
    \begin{equation*}
        u_2
        =
        \frac{3.5\times10^{-6}\times10\,\mathrm{V}}{3}
        =
        1.17\times10^{-5}\,\mathrm{V}
        =
        11.7\,\mu\mathrm{V}
    \end{equation*}

    其相对标准差为 $25\%$，则自由度为
    \begin{equation*}
        \nu_2
        =
        \frac{1}{2\left(\dfrac{\sigma_{u_2}}{u_2}\right)^2}
        =
        \frac{1}{2\times(25\%)^2}
        =
        8
    \end{equation*}

    \item 标准电压表的稳定度引起的标准不确定度 $u_3$

    标准电压表的稳定度按均匀分布，其不确定度 $u_3$ 属 B 类评定，根据式（4-3）计算标准不确定度为
    \begin{equation*}
        u_3
        =
        \frac{15\,\mu\mathrm{V}}{\sqrt{3}}
        =
        8.7\,\mu\mathrm{V}
    \end{equation*}

    其相对标准差为 $10\%$，则自由度为
    \begin{equation*}
        \nu_3
        =
        \frac{1}{2\left(\dfrac{\sigma_{u_3}}{u_3}\right)^2}
        =
        \frac{1}{2\times(10\%)^2}
        =
        50
    \end{equation*}
\end{enumerate}

\problem{例4-4}

\timu 已知被测量的估计值为 $20.0005\,\mathrm{mm}$，若有两种情况：

\textcircled{1} 展伸不确定度 $U=0.00124\,\mathrm{mm}$；

\textcircled{2} 展伸不确定度 $U=0.00123\,\mathrm{mm}$。要求对 $U$ 进行修约。

\jieda 根据被测量的估计值，取 $0.0001\,\mathrm{mm}$ 作为 $U$ 的基本单位。

\textcircled{1} $U=0.00124\,\mathrm{mm}$，其整数部分为 $12$，小数部分为 $0.4$，大于基本单位的 $1/3$，故舍去后整数单位加 $1$。修约后，$U=0.0013\,\mathrm{mm}$；

\textcircled{2} $U=0.00123\,\mathrm{mm}$，其整数部分为 $12$，小数部分为 $0.3$，小于基本单位的 $1/3$，故舍去。修约后，$U=0.0012\,\mathrm{mm}$。

\problem{例4-5}

\timu 根据例4-3计算的各标准不确定度，并取置信概率 $P=95\%$，求电压测量的展伸不确定度及其自由度。

\jieda 在例4-3中，已计算出各误差源的标准不确定度 $u_i$；由于不确定度的传递系数 $\left|\dfrac{\partial f}{\partial x_i}\right|=1$，故各个不确定度分量与其标准不确定度相等，即 $u_i=u_{x_i}$。因不确定度分量 $u_1$、$u_2$、$u_3$ 相互独立，则不确定度的相关系数 $\rho_{ij}=0$。计算电压测量的合成标准不确定度为
\begin{equation*}
    u_c
    =
    \sqrt{u_1^2+u_2^2+u_3^2}
    =
    \sqrt{(2.8)^2+(11.7)^2+(8.7)^2}\,\mu\mathrm{V}
    =
    15\,\mu\mathrm{V}
\end{equation*}

计算自由度得
\begin{equation*}
    \nu
    =
    \frac{u_c^4}
    {
    \dfrac{u_1^4}{\nu_1}
    +
    \dfrac{u_2^4}{\nu_2}
    +
    \dfrac{u_3^4}{\nu_3}
    }
    =
    \frac{(15)^4}
    {
    \dfrac{(2.8)^4}{9}
    +
    \dfrac{(11.7)^4}{8}
    +
    \dfrac{(8.7)^4}{50}
    }
    =
    20
\end{equation*}

根据置信概率 $P=95\%$，自由度 $\nu=20$，查 $t$ 分布表得 $t_{0.95}(20)=2.09$，即包含因子 $k=2.09$。于是，由电测量的展伸不确定度为
\begin{equation*}
    U
    =
    ku_c
    =
    2.09\times15\,\mu\mathrm{V}
    =
    31.35\,\mu\mathrm{V}
\end{equation*}

依据“三分之一准则”对展伸不确定度进行修约，得展伸不确定度 $U=32\,\mu\mathrm{V}$。

最后的不确定度报告：电压测量结果为
\begin{equation*}
    V=(10.000104\pm0.000032)\,\mathrm{V}
\end{equation*}

说明：以上测量结果中 $\pm$ 符号后的数值是展伸不确定度 $U=32\,\mu\mathrm{V}$，是由合成标准不确定度 $u_c=0.000015\,\mathrm{V}$ 及包含因子 $k=2.09$ 确定的。对应的置信概率 $P=95\%$，自由度 $\nu=20$。


\section*{第五章\quad 线性测量的参数最小二乘法处理}

\problem{例5-1}

\timu 已知任意温度 $t$ 时的铜棒长度 $y_t$、$0^\circ\mathrm{C}$ 时的铜棒长度 $y_0$ 和铜的线膨胀系数 $\alpha$ 具有线性关系
\begin{equation*}
    y_t=y_0(1+\alpha t)
\end{equation*}
现测得在不同温度 $t_i$ 下铜棒长度 $l_i$ 如下表，试估计 $y_0$ 和 $\alpha$ 的最可信赖值。

\begin{center}
\small
\setlength{\tabcolsep}{5pt}
\begin{tabular}{c|cccccc}
\toprule
$i$ & 1 & 2 & 3 & 4 & 5 & 6\\
\midrule
$t_i/{}^\circ\mathrm{C}$ & 10 & 20 & 25 & 30 & 40 & 45\\
$l_i/\mathrm{mm}$ & 2000.36 & 2000.72 & 2000.80 & 2001.07 & 2001.48 & 2001.60\\
\bottomrule
\end{tabular}
\end{center}

\jieda 列出误差方程为
\begin{equation*}
    v_i=l_i-y_0(1+\alpha t_i),
    \qquad i=1,2,\cdots,6
\end{equation*}
其中，$l_i$ 为在温度 $t_i$ 下铜棒长度的测得值，$\alpha$ 为铜的线膨胀系数。

令
\begin{equation*}
    y_0=a,\qquad \alpha y_0=b
\end{equation*}
则误差方程可写为
\begin{equation*}
    v_i=l_i-(a+t_i b),
    \qquad i=1,2,\cdots,6
\end{equation*}

为计算方便，将数据列表如下：

\begin{center}
\small
\setlength{\tabcolsep}{5pt}
\begin{tabular}{c|c|c|c|c}
\toprule
$i$ & $t_i/{}^\circ\mathrm{C}$ & $t_i^2/{}^\circ\mathrm{C}^2$ & $l_i/\mathrm{mm}$ & $t_i l_i/({}^\circ\mathrm{C}\cdot\mathrm{mm})$\\
\midrule
1 & 10 & 100 & 2000.36 & 20003.6\\
2 & 20 & 400 & 2000.72 & 40014.4\\
3 & 25 & 625 & 2000.80 & 50020.0\\
4 & 30 & 900 & 2001.07 & 60032.1\\
5 & 40 & 1600 & 2001.48 & 80059.2\\
6 & 45 & 2025 & 2001.60 & 90072.0\\
\midrule
$\sum$ & 170 & 5650 & 12006.03 & 340201.3\\
\bottomrule
\end{tabular}
\end{center}

根据误差方程，列出正规方程为
\begin{equation*}
\left\{
\begin{aligned}
    na+\sum_{i=1}^{6}t_i b&=\sum_{i=1}^{6}l_i\\
    \sum_{i=1}^{6}t_i a+\sum_{i=1}^{6}t_i^2 b&=\sum_{i=1}^{6}t_i l_i
\end{aligned}
\right.
\end{equation*}

将表中数据代入，得
\begin{equation*}
\left\{
\begin{aligned}
    6a+170b&=12006.03\\
    170a+5650b&=340201.3
\end{aligned}
\right.
\end{equation*}

解得
\begin{equation*}
    a\approx1999.97\,\mathrm{mm},
    \qquad
    b\approx0.03654\,\mathrm{mm}/{}^\circ\mathrm{C}
\end{equation*}

即
\begin{equation*}
    y_0=a=1999.97\,\mathrm{mm}
\end{equation*}
\begin{equation*}
    \alpha=\frac{b}{y_0}
    =
    \frac{0.03654\,\mathrm{mm}/{}^\circ\mathrm{C}}{1999.97\,\mathrm{mm}}
    \approx1.83\times10^{-5}/{}^\circ\mathrm{C}
\end{equation*}

按矩阵形式求解，则有
\begin{equation*}
    C=
    \begin{pmatrix}
        n & \displaystyle\sum_{i=1}^{6}t_i\\
        \displaystyle\sum_{i=1}^{6}t_i & \displaystyle\sum_{i=1}^{6}t_i^2
    \end{pmatrix}
    =
    \begin{pmatrix}
        6 & 170\\
        170 & 5650
    \end{pmatrix}
,\quad
    C^{-1}\approx
    \begin{pmatrix}
        1.13 & -0.034\\
        -0.034 & 0.0012
    \end{pmatrix}
\end{equation*}
\begin{equation*}
    A^{\mathrm{T}}L=
    \begin{pmatrix}
        \displaystyle\sum_{i=1}^{6}l_i\\
        \displaystyle\sum_{i=1}^{6}t_i l_i
    \end{pmatrix}
    =
    \begin{pmatrix}
        12006.03\\
        340201.3
    \end{pmatrix}
\end{equation*}

因此
\begin{equation*}
    \hat{X}
    =
    \begin{pmatrix}
        a\\
        b
    \end{pmatrix}
    =
    C^{-1}A^{\mathrm{T}}L
    =
    \begin{pmatrix}
        1999.97\\
        0.03654
    \end{pmatrix}
\end{equation*}

所以
\begin{equation*}
    y_0=1999.97\,\mathrm{mm},
    \qquad
    \alpha\approx1.83\times10^{-5}/{}^\circ\mathrm{C}
\end{equation*}

因此，铜棒长度 $y_t$ 随温度 $t$ 的线性变化规律为
\begin{equation*}
    y_t=1999.97\left(1+1.83\times10^{-5}t\right)\,\mathrm{mm}
\end{equation*}

\problem{例5-2}

\timu 某测量过程中误差方程式及相应的标准差如下：
\begin{equation*}
\left\{
\begin{aligned}
    v_1&=6.44-(x_1+x_2),        & \sigma_1&=0.06\\
    v_2&=8.60-(x_1+2x_2),       & \sigma_2&=0.06\\
    v_3&=10.81-(x_1+3x_2),      & \sigma_3&=0.08\\
    v_4&=13.22-(x_1+4x_2),      & \sigma_4&=0.08\\
    v_5&=15.27-(x_1+5x_2),      & \sigma_5&=0.08
\end{aligned}
\right.
\end{equation*}
试求 $x_1$、$x_2$ 的最小二乘法处理正规方程的解。

\jieda 首先确定各式的权，由权与方差成反比可得
\begin{equation*}
    p_1:p_2:p_3:p_4:p_5
    =
    \frac{1}{\sigma_1^2}:
    \frac{1}{\sigma_2^2}:
    \frac{1}{\sigma_3^2}:
    \frac{1}{\sigma_4^2}:
    \frac{1}{\sigma_5^2}
    =
    16:16:9:9:9
\end{equation*}

取各式的权为
\begin{equation*}
    p_1=16,\qquad p_2=16,\qquad p_3=9,\qquad p_4=9,\qquad p_5=9
\end{equation*}

现用表格计算给出正规方程常数项和系数项：

\begin{center}
\scriptsize
\setlength{\tabcolsep}{4pt}
\begin{tabular}{c|c|c|c|c|c|c|c|c|c}
\toprule
$i$ & $a_1$ & $a_2$ & $p_i$ & $p_i a_1^2$ & $p_i a_2^2$ & $p_i a_1a_2$ & $l_i$ & $p_i a_1l_i$ & $p_i a_2l_i$\\
\midrule
1 & 1 & 1 & 16 & 16 & 16 & 16 & 6.44 & 103.04 & 103.04\\
2 & 1 & 2 & 16 & 16 & 64 & 32 & 8.60 & 137.60 & 275.20\\
3 & 1 & 3 & 9 & 9 & 81 & 27 & 10.81 & 97.29 & 291.87\\
4 & 1 & 4 & 9 & 9 & 144 & 36 & 13.22 & 118.98 & 475.92\\
5 & 1 & 5 & 9 & 9 & 225 & 45 & 15.27 & 137.43 & 687.15\\
\midrule
$\sum$ &  &  &  & 59 & 530 & 156 &  & 594.34 & 1833.18\\
\bottomrule
\end{tabular}
\end{center}

可得正规方程为
\begin{equation*}
\left\{
\begin{aligned}
    59x_1+156x_2&=594.34\\
    156x_1+530x_2&=1833.18
\end{aligned}
\right.
\end{equation*}

解得最小二乘法处理结果为
\begin{equation*}
    x_1\approx4.186,
    \qquad
    x_2\approx2.227
\end{equation*}

\problem{例5-3}

\timu 试求例5-1中铜棒长度的测量精度。

\jieda 已知残余误差方程为
\begin{equation*}
    v_i=
    \left[
    l_i-1999.97\left(1+1.83\times10^{-5}t_i\right)
    \right]\mathrm{mm},
    \qquad i=1,2,\cdots,6
\end{equation*}

将 $t_i$、$l_i$ 值代入上式，可得残余误差为
\begin{equation*}
\begin{aligned}
    v_1&=
    \left[
    2000.36-1999.97\left(1+1.83\times10^{-5}\times10\right)
    \right]\mathrm{mm}
    \approx0.03\,\mathrm{mm}\\
    v_2&=
    \left[
    2000.72-1999.97\left(1+1.83\times10^{-5}\times20\right)
    \right]\mathrm{mm}
    \approx0.02\,\mathrm{mm}\\
    v_3&=
    \left[
    2000.80-1999.97\left(1+1.83\times10^{-5}\times25\right)
    \right]\mathrm{mm}
    \approx-0.08\,\mathrm{mm}\\
    v_4&=
    \left[
    2001.07-1999.97\left(1+1.83\times10^{-5}\times30\right)
    \right]\mathrm{mm}
    \approx0\,\mathrm{mm}\\
    v_5&=
    \left[
    2001.48-1999.97\left(1+1.83\times10^{-5}\times40\right)
    \right]\mathrm{mm}
    \approx0.05\,\mathrm{mm}\\
    v_6&=
    \left[
    2001.60-1999.97\left(1+1.83\times10^{-5}\times45\right)
    \right]\mathrm{mm}
    \approx-0.02\,\mathrm{mm}
\end{aligned}
\end{equation*}

于是
\begin{equation*}
    \sum_{i=1}^{6}v_i^2=0.0106\,\mathrm{mm}^2
\end{equation*}

因 $n=6$，未知参数个数为 $t=2$，所以可得标准差为
\begin{equation*}
    \sigma
    =
    \sqrt{
    \frac{\displaystyle\sum_{i=1}^{6}v_i^2}{n-t}
    }
    =
    \sqrt{\frac{0.0106}{6-2}}\,\mathrm{mm}
    =
    0.051\,\mathrm{mm}
\end{equation*}

\problem{例5-4}

\timu 试求例5-1中铜棒长度和线膨胀系数估计量的精度。

\jieda 已知正规方程为
\begin{equation*}
\left\{
\begin{aligned}
    6a+170b&=12006.03\\
    170a+5650b&=340201.3
\end{aligned}
\right.
\end{equation*}

测量数据 $l_i$ 的标准差为
\begin{equation*}
    \sigma=0.051\,\mathrm{mm}
\end{equation*}

由正规方程系数可列出求解不定乘数的方程组。设
\begin{equation*}
    C^{-1}=
    \begin{pmatrix}
        d_{11} & d_{12}\\
        d_{21} & d_{22}
    \end{pmatrix}
\end{equation*}
则对应方程为
\begin{equation*}
\left\{
\begin{aligned}
    6d_{11}+170d_{12}&=1\\
    170d_{11}+5650d_{12}&=0\\
    6d_{21}+170d_{22}&=0\\
    170d_{21}+5650d_{22}&=1
\end{aligned}
\right.
\end{equation*}

分别解得
\begin{equation*}
    d_{11}=1.13,
    \qquad
    d_{22}=0.0012
\end{equation*}

则由参数标准差公式，可得估计量 $a$、$b$ 的标准差为
\begin{equation*}
    \sigma_a
    =
    \sigma\sqrt{d_{11}}
    =
    0.051\sqrt{1.13}\,\mathrm{mm}
    =
    0.054\,\mathrm{mm}
\end{equation*}
\begin{equation*}
    \sigma_b
    =
    \sigma\sqrt{d_{22}}
    =
    0.051\sqrt{0.0012}\,\mathrm{mm}/{}^\circ\mathrm{C}
    =
    0.0018\,\mathrm{mm}/{}^\circ\mathrm{C}
\end{equation*}

因
\begin{equation*}
    y_0=a,
    \qquad
    \alpha=\frac{b}{y_0}
\end{equation*}

故
\begin{equation*}
    \sigma_{y_0}
    =
    \sigma_a
    =
    0.054\,\mathrm{mm}
\end{equation*}
\begin{equation*}
    \sigma_{\alpha}
    =
    \frac{\sigma_b}{y_0}
    =
    \frac{0.0018}{1999.97}/{}^\circ\mathrm{C}
    \approx9\times10^{-7}/{}^\circ\mathrm{C}
\end{equation*}

\problem{例5-5}

\timu 为了测定公称值是 $10\,\mathrm{g}$、$20\,\mathrm{g}$、$50\,\mathrm{g}$ 的三只砝码质量 $m_1$、$m_2$、$m_3$，对三只砝码采用不同组合进行了 $7$ 次测量。测量无系统误差，按下列组合方式和测量后的结果为：
\begin{equation*}
\begin{aligned}
    \text{对砝码 }m_1\text{ 单独测量，得}\quad &y_1=10.002\,\mathrm{g}\\
    \text{对砝码 }m_2\text{ 单独测量，得}\quad &y_2=20.002\,\mathrm{g}\\
    \text{对砝码 }m_3\text{ 单独测量，得}\quad &y_3=50.006\,\mathrm{g}\\
    \text{对砝码 }m_1,m_2\text{ 之和组合测量，得}\quad &y_4=30.004\,\mathrm{g}\\
    \text{对砝码 }m_1,m_3\text{ 之和组合测量，得}\quad &y_5=60.002\,\mathrm{g}\\
    \text{对砝码 }m_2,m_3\text{ 之和组合测量，得}\quad &y_6=70.002\,\mathrm{g}\\
    \text{对砝码 }m_1,m_2,m_3\text{ 之和组合测量，得}\quad &y_7=80.008\,\mathrm{g}
\end{aligned}
\end{equation*}
求三只砝码的最佳估计值及其标准差。

\jieda 首先列出误差方程：
\begin{equation*}
\left\{
\begin{aligned}
    v_1&=y_1-m_1\\
    v_2&=y_2-m_2\\
    v_3&=y_3-m_3\\
    v_4&=y_4-(m_1+m_2)\\
    v_5&=y_5-(m_1+m_3)\\
    v_6&=y_6-(m_2+m_3)\\
    v_7&=y_7-(m_1+m_2+m_3)
\end{aligned}
\right.
\end{equation*}

根据矩阵形式，上式可以写为
\begin{equation*}
    V=Y-A\hat{M}
\end{equation*}
其中
\begin{equation*}
    V=
    \begin{pmatrix}
        v_1\\
        v_2\\
        v_3\\
        v_4\\
        v_5\\
        v_6\\
        v_7
    \end{pmatrix},
    \qquad
    Y=
    \begin{pmatrix}
        10.002\\
        20.002\\
        50.006\\
        30.004\\
        60.002\\
        70.002\\
        80.008
    \end{pmatrix},
    \qquad
    A=
    \begin{pmatrix}
        1 & 0 & 0\\
        0 & 1 & 0\\
        0 & 0 & 1\\
        1 & 1 & 0\\
        1 & 0 & 1\\
        0 & 1 & 1\\
        1 & 1 & 1
    \end{pmatrix},
    \qquad
    \hat{M}=
    \begin{pmatrix}
        m_1\\
        m_2\\
        m_3
    \end{pmatrix}
\end{equation*}

由式
\begin{equation*}
    \hat{M}=C^{-1}A^{\mathrm{T}}Y
\end{equation*}
可得
\begin{equation*}
    C^{-1}
    =
    (A^{\mathrm{T}}A)^{-1}
    =
    \begin{pmatrix}
        4 & 2 & 2\\
        2 & 4 & 2\\
        2 & 2 & 4
    \end{pmatrix}^{-1}
    =
    \frac{1}{8}
    \begin{pmatrix}
        3 & -1 & -1\\
        -1 & 3 & -1\\
        -1 & -1 & 3
    \end{pmatrix}
\end{equation*}

所以
\begin{equation*}
\begin{aligned}
    \hat{M}
    &=
    \begin{pmatrix}
        m_1\\
        m_2\\
        m_3
    \end{pmatrix}
    =
    C^{-1}A^{\mathrm{T}}Y\\
    &=
    \frac{1}{8}
    \begin{pmatrix}
        3 & -1 & -1\\
        -1 & 3 & -1\\
        -1 & -1 & 3
    \end{pmatrix}
    \begin{pmatrix}
        1 & 0 & 0 & 1 & 1 & 0 & 1\\
        0 & 1 & 0 & 1 & 0 & 1 & 1\\
        0 & 0 & 1 & 0 & 1 & 1 & 1
    \end{pmatrix}
    \begin{pmatrix}
        10.002\\
        20.002\\
        50.006\\
        30.004\\
        60.002\\
        70.002\\
        80.008
    \end{pmatrix}
\end{aligned}
\end{equation*}

即
\begin{equation*}
    \hat{M}
    =
    \frac{1}{8}
    \begin{pmatrix}
        80.014\\
        160.014\\
        400.022
    \end{pmatrix}
    =
    \begin{pmatrix}
        10.00175\\
        20.00175\\
        50.00275
    \end{pmatrix}
\end{equation*}

即解得
\begin{equation*}
    m_1=10.00175\,\mathrm{g},
    \qquad
    m_2=20.00175\,\mathrm{g},
    \qquad
    m_3=50.00275\,\mathrm{g}
\end{equation*}

这就是砝码质量 $m_1$、$m_2$、$m_3$ 的最佳估计值，再求其上述估计量的标准差。

将最佳估计值代入误差方程可得
\begin{equation*}
\begin{aligned}
    v_1&=y_1-m_1=10.002-10.00175=0.00025\,\mathrm{g}\\
    v_2&=y_2-m_2=20.002-20.00175=0.00025\,\mathrm{g}\\
    v_3&=y_3-m_3=50.006-50.00275=0.00325\,\mathrm{g}\\
    v_4&=y_4-(m_1+m_2)=30.004-(10.00175+20.00175)=0.0005\,\mathrm{g}\\
    v_5&=y_5-(m_1+m_3)=60.002-(10.00175+50.00275)=-0.0025\,\mathrm{g}\\
    v_6&=y_6-(m_2+m_3)=70.002-(20.00175+50.00275)=-0.0025\,\mathrm{g}\\
    v_7&=y_7-(m_1+m_2+m_3)=80.008-(10.00175+20.00175+50.00275)=0.00175\,\mathrm{g}
\end{aligned}
\end{equation*}

于是
\begin{equation*}
\begin{aligned}
    \sum_{i=1}^{7}v_i^2
    &=
    0.00025^2+0.00025^2+0.00325^2+0.0005^2\\
    &\quad
    +(-0.0025)^2+(-0.0025)^2+0.00175^2\\
    &=
    0.0000265\,\mathrm{g}^2
\end{aligned}
\end{equation*}

因为是等精度测量，测得数据 $y_1,y_2,\cdots,y_7$ 的标准差相同，为
\begin{equation*}
    \sigma
    =
    \sqrt{
    \frac{\displaystyle\sum_{i=1}^{7}v_i^2}{n-t}
    }
    =
    \sqrt{\frac{0.0000265}{7-3}}\,\mathrm{g}
    =
    0.0026\,\mathrm{g}
\end{equation*}

为求估计量 $m_1,m_2,m_3$ 的标准差，首先求出不定乘数 $d_{ij}$。由式可知，不定乘数 $d_{ij}$ 的系数与正规方程的系数相同，因而 $d_{ij}$ 是矩阵 $C^{-1}$ 中的各元素，即
\begin{equation*}
    C^{-1}
    =
    \begin{pmatrix}
        d_{11} & d_{12} & d_{13}\\
        d_{21} & d_{22} & d_{23}\\
        d_{31} & d_{32} & d_{33}
    \end{pmatrix}
    =
    \frac{1}{8}
    \begin{pmatrix}
        3 & -1 & -1\\
        -1 & 3 & -1\\
        -1 & -1 & 3
    \end{pmatrix}
\end{equation*}

则
\begin{equation*}
    d_{11}=\frac{3}{8}=0.375
\end{equation*}

\begin{equation*}
    d_{22}=\frac{3}{8}=0.375
\end{equation*}

\begin{equation*}
    d_{33}=\frac{3}{8}=0.375
\end{equation*}

按照参数估计量标准差公式，可得
\begin{equation*}
    \sigma_{m_1}
    =
    \sigma\sqrt{d_{11}}
    =
    0.0026\sqrt{0.375}\,\mathrm{g}
    =
    0.0016\,\mathrm{g}
\end{equation*}
\begin{equation*}
    \sigma_{m_2}
    =
    \sigma\sqrt{d_{22}}
    =
    0.0026\sqrt{0.375}\,\mathrm{g}
    =
    0.0016\,\mathrm{g}
\end{equation*}
\begin{equation*}
    \sigma_{m_3}
    =
    \sigma\sqrt{d_{33}}
    =
    0.0026\sqrt{0.375}\,\mathrm{g}
    =
    0.0016\,\mathrm{g}
\end{equation*}

所以三只砝码的最佳估计值及其标准差为
\begin{equation*}
    m_1=10.00175\,\mathrm{g},
    \qquad
    \sigma_{m_1}=0.0016\,\mathrm{g}
\end{equation*}
\begin{equation*}
    m_2=20.00175\,\mathrm{g},
    \qquad
    \sigma_{m_2}=0.0016\,\mathrm{g}
\end{equation*}
\begin{equation*}
    m_3=50.00275\,\mathrm{g},
    \qquad
    \sigma_{m_3}=0.0016\,\mathrm{g}
\end{equation*}


\section*{第六章\quad 回归分析}

\problem{例6-1}

\timu 测量某导线在一定温度 $x$ 下的电阻值 $y$ 得如下结果：

\begin{center}
\small
\setlength{\tabcolsep}{5pt}
\begin{tabular}{c|ccccccc}
\toprule
$x/{}^\circ\mathrm{C}$ & 19.1 & 25.0 & 30.1 & 36.0 & 40.0 & 46.5 & 50.0\\
\midrule
$y/\Omega$ & 76.30 & 77.80 & 79.75 & 80.80 & 82.35 & 83.90 & 85.10\\
\bottomrule
\end{tabular}
\end{center}

试找出它们之间的内在关系。

\jieda 由一元线性回归方程的计算公式可知，回归方程的具体计算通常可以通过列表进行。计算表如下：

\begin{center}
\small
\setlength{\tabcolsep}{4.5pt}
\begin{tabular}{c|c|c|c|c|c}
\toprule
序号 & $x_i/{}^\circ\mathrm{C}$ & $y_i/\Omega$ & $x_i^2/{}^\circ\mathrm{C}^2$ & $y_i^2/\Omega^2$ & $x_i y_i/(\Omega\cdot{}^\circ\mathrm{C})$\\
\midrule
1 & 19.1 & 76.30 & 364.81 & 5821.690 & 1457.330\\
2 & 25.0 & 77.80 & 625.00 & 6052.840 & 1945.000\\
3 & 30.1 & 79.75 & 906.01 & 6360.062 & 2400.475\\
4 & 36.0 & 80.80 & 1296.00 & 6528.640 & 2908.800\\
5 & 40.0 & 82.35 & 1600.00 & 6781.522 & 3294.000\\
6 & 46.5 & 83.90 & 2162.25 & 7039.210 & 3901.350\\
7 & 50.0 & 85.10 & 2500.00 & 7242.010 & 4255.000\\
\midrule
$\sum$ & 246.7 & 566.00 & 9454.07 & 45825.974 & 20161.955\\
\bottomrule
\end{tabular}
\end{center}

由表中数据可得
\begin{equation*}
    \sum_{i=1}^{N}x_i=246.7{}^\circ\mathrm{C}, \qquad \bar{x}=35.243{}^\circ\mathrm{C}, \qquad \sum_{i=1}^{N}x_i^2=9454.07{}^\circ\mathrm{C}^2
\end{equation*}

\begin{equation*}
    \frac{\left(\displaystyle\sum_{i=1}^{N}x_i\right)^2}{N} = 8694.413{}^\circ\mathrm{C}^2
\end{equation*}

于是
\begin{equation*}
    l_{xx} = \sum_{i=1}^{N}x_i^2 - \frac{\left(\displaystyle\sum_{i=1}^{N}x_i\right)^2}{N} = 759.657{}^\circ\mathrm{C}^2
\end{equation*}

同理有 $\sum_{i=1}^{N}y_i=566.00\Omega, \qquad \bar{y}=80.857\Omega, \qquad \sum_{i=1}^{N}y_i^2=45825.974\Omega^2$

\begin{equation*}
    \frac{\left(\displaystyle\sum_{i=1}^{N}y_i\right)^2}{N} = 45765.143\Omega^2
\end{equation*}

于是
\begin{equation*}
    l_{yy} = \sum_{i=1}^{N}y_i^2 - \frac{\left(\displaystyle\sum_{i=1}^{N}y_i\right)^2}{N} = 60.831\Omega^2
\end{equation*}

又因为 $\sum_{i=1}^{N}x_i y_i=20161.955\Omega\cdot{}^\circ\mathrm{C}$

\begin{equation*}
    \frac{\left(\displaystyle\sum_{i=1}^{N}x_i\right) \left(\displaystyle\sum_{i=1}^{N}y_i\right)}{N} = 19947.457\Omega\cdot{}^\circ\mathrm{C}
\end{equation*}

所以
\begin{equation*}
    l_{xy} = \sum_{i=1}^{N}x_i y_i - \frac{\left(\displaystyle\sum_{i=1}^{N}x_i\right) \left(\displaystyle\sum_{i=1}^{N}y_i\right)}{N} = 214.498\Omega\cdot{}^\circ\mathrm{C}
\end{equation*}

因此回归系数为 $b=l_{xy}/l_{xx}=214.498/759.657=0.2824\Omega/{}^\circ\mathrm{C}$，回归常数为 $b_0=\bar{y}-b\bar{x}=80.857-0.2824\times35.243=70.90\Omega$

故所求回归方程为 $\hat{y} = 70.90\Omega + (0.2824\Omega/{}^\circ\mathrm{C})x$

这条回归直线一定通过 $(\bar{x},\bar{y})$ 这一点。再令 $x=x_0$，代入回归方程，可求得相应的 $y_0$，连接 $(\bar{x},\bar{y})$ 和 $(x_0,y_0)$，即为所求回归直线。在本例中，回归系数 $b$ 的物理意义是温度每上升 $1{}^\circ\mathrm{C}$，电阻平均增加 $0.2824\Omega$。

\problem{例6-2}

\timu 在例6-1中，对电阻和温度的回归方程进行方差分析，判断回归方程是否显著。

\jieda 把平方和及自由度进行分解的方法，可以归纳为方差分析表。对于一元线性回归，有如下形式：

\begin{center}
\small
\setlength{\tabcolsep}{5pt}
\begin{tabular}{c|c|c|c|c|c}
\toprule
来源 & 平方和 & 自由度 & 方差 & $F$ & 显著性\\
\midrule
回归 & $U=bl_{xy}$ & $1$ & -- & $\dfrac{U/1}{Q/(N-2)}$ & --\\
残余 & $Q=l_{yy}-bl_{xy}$ & $N-2$ & $\sigma^2=\dfrac{Q}{N-2}$ & -- & --\\
总计 & $S=l_{yy}$ & $N-1$ & -- & -- & --\\
\bottomrule
\end{tabular}
\end{center}

由例6-1的计算结果可得 $U=bl_{xy}=0.2824\times214.498=60.574\Omega^2$，$Q=l_{yy}-U=60.831-60.574=0.257\Omega^2$，$\sigma^2=Q/(N-2)=0.257/(7-2)=0.0514\Omega^2$

所以方差分析表为

\begin{center}
\small
\setlength{\tabcolsep}{5pt}
\begin{tabular}{c|c|c|c|c|c}
\toprule
来源 & 平方和$/\Omega^2$ & 自由度 & 方差$/\Omega^2$ & $F$ & 显著性\\
\midrule
回归 & 60.574 & 1 & 60.574 & $1.18\times10^3$ & $\alpha=0.01$\\
残余 & 0.257 & 5 & 0.0514 & -- & --\\
总计 & 60.831 & 6 & -- & -- & --\\
\bottomrule
\end{tabular}
\end{center}

显著性一栏中的 $\alpha=0.01$，表明前面所得回归方程在 $\alpha=0.01$ 水平上显著，即可信赖程度为 $99\%$ 以上，这是高度显著的。

\problem{例6-3}

\timu 用标准压力计对某固体压力传感器进行检定，检定所得数据见下表。表中 $x_i$ 为标准压力，$y_i$ 为传感器输出电压，$\bar{y}_i$ 为 $4$ 次读数的算术平均值。试对仪器定标，并分析仪器的误差。

\begin{center}
\scriptsize
\setlength{\tabcolsep}{4pt}
\begin{tabular}{c|c|cccc|c}
\toprule
序号 & $x_i/(\mathrm{N}\cdot\mathrm{cm}^{-2})$ & \multicolumn{4}{c|}{$y_i/\mathrm{mV}$} & $\bar{y}_i/\mathrm{mV}$\\
\cline{3-6}
 &  & 升压 & 降压 & 升压 & 降压 & \\
\midrule
1 & 0  & 2.78 & 2.80 & 2.80 & 2.86 & 2.8100\\
2 & 1  & 9.70 & 9.76 & 9.78 & 9.78 & 9.7550\\
3 & 2  & 16.60 & 16.71 & 16.70 & 16.76 & 16.6925\\
4 & 3  & 23.54 & 23.56 & 23.58 & 23.71 & 23.5975\\
5 & 4  & 30.44 & 30.51 & 30.54 & 30.64 & 30.5325\\
6 & 5  & 37.35 & 37.45 & 37.42 & 37.50 & 37.4300\\
7 & 6  & 44.28 & 44.35 & 44.30 & 44.38 & 44.3275\\
8 & 7  & 51.19 & 51.25 & 51.18 & 51.25 & 51.2175\\
9 & 8  & 58.06 & 58.08 & 58.12 & 58.14 & 58.1000\\
10 & 9 & 64.92 & 64.96 & 64.94 & 65.00 & 64.9550\\
11 & 10 & 71.73 & 71.73 & 71.75 & 71.75 & 71.7400\\
\bottomrule
\end{tabular}
\end{center}

\jieda 要求线性定标，故应拟合一条回归直线。可以证明，用平均值的 $11$ 个点拟合的回归直线与用原来的 $44$ 个点拟合的回归直线完全一样。具体计算如下，其中 $y_i$ 即为上表中的 $\bar{y}_i$。

\begin{center}
\scriptsize
\setlength{\tabcolsep}{5pt}
\begin{tabular}{c|c|c|c|c|c}
\toprule
序号 & $x_i/(\mathrm{N}\cdot\mathrm{cm}^{-2})$ & $y_i/\mathrm{mV}$ & $x_i^2/(\mathrm{N}\cdot\mathrm{cm}^{-2})^2$ & $y_i^2/\mathrm{mV}^2$ & $x_i y_i/(\mathrm{mV}\cdot\mathrm{N}\cdot\mathrm{cm}^{-2})$\\
\midrule
1 & 0  & 2.8100 & 0 & 7.8961 & 0\\
2 & 1  & 9.7550 & 1 & 95.1600 & 9.7550\\
3 & 2  & 16.6925 & 4 & 278.6396 & 33.3850\\
4 & 3  & 23.5975 & 9 & 556.8420 & 70.7925\\
5 & 4  & 30.5325 & 16 & 932.2336 & 122.1300\\
6 & 5  & 37.4300 & 25 & 1401.0049 & 187.1500\\
7 & 6  & 44.3275 & 36 & 1964.9273 & 265.9650\\
8 & 7  & 51.2175 & 49 & 2623.2323 & 358.5225\\
9 & 8  & 58.1000 & 64 & 3375.6100 & 464.8000\\
10 & 9 & 64.9550 & 81 & 4219.1520 & 584.5950\\
11 & 10 & 71.7400 & 100 & 5146.6276 & 717.4000\\
\midrule
$\sum$ & 55 & 411.1575 & 385 & 20601.3254 & 2814.4950\\
\bottomrule
\end{tabular}
\end{center}

由表中数据可得
\begin{equation*}
    \sum_{i=1}^{N}x_i=55\mathrm{N}\cdot\mathrm{cm}^{-2}, \qquad \bar{x}=5\mathrm{N}\cdot\mathrm{cm}^{-2}, \qquad \sum_{i=1}^{N}x_i^2=385(\mathrm{N}\cdot\mathrm{cm}^{-2})^2
\end{equation*}

\begin{equation*}
    \frac{\left(\displaystyle\sum_{i=1}^{N}x_i\right)^2}{N} = 275(\mathrm{N}\cdot\mathrm{cm}^{-2})^2
\end{equation*}

所以
\begin{equation*}
    l_{xx} = \sum_{i=1}^{N}x_i^2 - \frac{\left(\displaystyle\sum_{i=1}^{N}x_i\right)^2}{N} = 110(\mathrm{N}\cdot\mathrm{cm}^{-2})^2
\end{equation*}

又有
\begin{equation*}
    \sum_{i=1}^{N}y_i=411.1575\mathrm{mV}, \qquad \bar{y}=37.3780\mathrm{mV}, \qquad \sum_{i=1}^{N}y_i^2=20601.3254\mathrm{mV}^2
\end{equation*}

\begin{equation*}
    \frac{\left(\displaystyle\sum_{i=1}^{N}y_i\right)^2}{N} = 15368.226\mathrm{mV}^2
\end{equation*}

于是
\begin{equation*}
    l_{yy} = \sum_{i=1}^{N}y_i^2 - \frac{\left(\displaystyle\sum_{i=1}^{N}y_i\right)^2}{N} = 5233.0990\mathrm{mV}^2
\end{equation*}

又因为 $\sum_{i=1}^{N}x_i y_i=2814.4950\mathrm{mV}\cdot\mathrm{N}\cdot\mathrm{cm}^{-2}$

\begin{equation*}
    \frac{\left(\displaystyle\sum_{i=1}^{N}x_i\right) \left(\displaystyle\sum_{i=1}^{N}y_i\right)}{N} = 2055.7875\mathrm{mV}\cdot\mathrm{N}\cdot\mathrm{cm}^{-2}
\end{equation*}

所以
\begin{equation*}
    l_{xy} = \sum_{i=1}^{N}x_i y_i - \frac{\left(\displaystyle\sum_{i=1}^{N}x_i\right) \left(\displaystyle\sum_{i=1}^{N}y_i\right)}{N} = 758.7075\mathrm{mV}\cdot\mathrm{N}\cdot\mathrm{cm}^{-2}
\end{equation*}

因此 $b = \frac{l_{xy}}{l_{xx}} = \frac{758.7075}{110} = 6.8973\mathrm{mV}/(\mathrm{N}\cdot\mathrm{cm}^{-2})$

$b_0 = \bar{y}-b\bar{x} = 37.3780-6.8973\times5 = 2.8913\mathrm{mV}$

故回归方程为 $\hat{y} = 2.8913\mathrm{mV} + \left[ 6.8973\mathrm{mV}/(\mathrm{N}\cdot\mathrm{cm}^{-2}) \right]x$

现在进行方差分析。当用 $\bar{y}_i$ 求回归直线时，各平方和可按下列顺序计算：
\begin{equation*}
\left\{
\begin{aligned}
    U&=mbl_{xy}\\
    Q_L&=ml_{yy}-U\\
    Q_E&=\sum_{i=1}^{N}\sum_{j=1}^{m}(y_{ij}-\bar{y}_i)^2\\
    S&=U+Q_L+Q_E
\end{aligned}
\right.
\end{equation*}

计算结果如下表所示，其中 $F$ 栏为用误差平方和和相应的自由度对各项进行 $F$ 检验的数学统计量。

\begin{center}
\small
\setlength{\tabcolsep}{5pt}
\begin{tabular}{c|c|c|c|c|c}
\toprule
来源 & 平方和$/\mathrm{mV}^2$ & 自由度 & 方差$/\mathrm{mV}^2$ & $F$ & 显著性\\
\midrule
回归 & 20932.2574 & 1 & 20932.2574 & $7.68\times10^6$ & $F_{0.01}=7.47$\\
失拟 & 0.1386 & 9 & 0.0154 & 5.65 & $F_{0.01}=3.03$\\
误差 & 0.0899 & 33 & 0.0027 & -- & --\\
总计 & 20932.4859 & 43 & -- & -- & --\\
\bottomrule
\end{tabular}
\end{center}

第一行有 $F = \frac{U/\nu_U}{Q_E/\nu_{Q_E}} = 7.68\times10^6$

由于 $F>F_{0.01}(1,33)=7.47$ 故回归高度显著。

第二行有 $F_1 = \frac{Q_L/\nu_{Q_L}}{Q_E/\nu_{Q_E}} = 5.65$

由于 $F_1>F_{0.01}(9,33)=3.03$ 故失拟平方和的检验结果也高度显著，说明失拟误差相对于实验误差来说是不可忽略的。这时有如下几种可能：
\begin{enumerate}[label=\arabic*)]
    \item 影响 $y$ 的除 $x$ 外，至少还有一个不可忽略的因素；
    \item $y$ 和 $x$ 是曲线关系；
    \item $y$ 和 $x$ 无关。
\end{enumerate}

\problem{例6-4}

\timu 对例6-1用分组法求回归方程。

\jieda 因观测数据已按自变量从小到大的次序排列，故可按实验顺序分成两组，并建立两组相应的观测方程。取 $k=(N+1)/2=4$，然后分别相加：

\begin{equation*}
\left\{
\begin{aligned}
    76.30\Omega&=b_0+(19.1{}^\circ\mathrm{C})b\\
    77.80\Omega&=b_0+(25.0{}^\circ\mathrm{C})b\\
    79.75\Omega&=b_0+(30.1{}^\circ\mathrm{C})b\\
    80.80\Omega&=b_0+(36.0{}^\circ\mathrm{C})b
\end{aligned}
\right.
\end{equation*}

相加得 $314.65\Omega=4b_0+(110.2{}^\circ\mathrm{C})b$。

另一组为
\begin{equation*}
\left\{
\begin{aligned}
    82.35\Omega&=b_0+(40.0{}^\circ\mathrm{C})b\\
    83.90\Omega&=b_0+(46.5{}^\circ\mathrm{C})b\\
    85.10\Omega&=b_0+(50.0{}^\circ\mathrm{C})b
\end{aligned}
\right.
\end{equation*}

相加得 $251.35\Omega=3b_0+(136.5{}^\circ\mathrm{C})b$。

解方程组
\begin{equation*}
\left\{
\begin{aligned}
    314.65\Omega&=4b_0+(110.2{}^\circ\mathrm{C})b\\
    251.35\Omega&=3b_0+(136.5{}^\circ\mathrm{C})b
\end{aligned}
\right.
\end{equation*}

得 $b_0=70.80\Omega, \qquad b=0.2853\Omega/{}^\circ\mathrm{C}$

故所求回归方程为 $\hat{y} = 70.80\Omega + (0.2853\Omega/{}^\circ\mathrm{C})x$

这与用最小二乘法求得的回归方程比较接近。此法简单明了，拟合的直线就是通过第一组重心和第二组重心的一条直线，是工程实践中常用的一种简单方法。

\problem{例6-5}

\timu 用 X 射线机检查镁合金焊接件及铸件内部缺陷时，为达到最佳灵敏度，透照电压 $y$ 应随被透照件厚度 $x$ 而改变。经实验得如下二组数据：

\begin{center}
\small
\setlength{\tabcolsep}{5pt}
\begin{tabular}{c|cccccccccc}
\toprule
$x/\mathrm{mm}$ & 12 & 13 & 15 & 16 & 18 & 20 & 22 & 24 & 25 & 26\\
\midrule
$y/\mathrm{kV}$ & 53 & 55 & 60 & 65 & 70 & 75 & 80 & 85 & 88 & 90\\
\bottomrule
\end{tabular}
\end{center}

\jieda 把这组数据点在坐标纸上，然后通过点群作一直线。在接近直线的两端各取一点，如 $(12,53)$ 和 $(25,88)$，即可求系数 $b = \frac{(88-53)\mathrm{kV}}{(25-12)\mathrm{mm}} = 2.7\mathrm{kV}/\mathrm{mm}$

将回归直线上任一点代入回归方程求出 $b_0$，如用 $53\mathrm{kV}=b_0+(2.7\mathrm{kV}/\mathrm{mm})\times12\mathrm{mm}$，得 $b_0=20.6\mathrm{kV}$

所以求得回归方程为 $\hat{y} = 20.6\mathrm{kV} + (2.7\mathrm{kV}/\mathrm{mm})x$

\problem{例6-6}

\timu 若例6-1中电阻值 $y$ 是可以精确测定的，而把所有误差都归结于温度 $x$，试求另一回归方程。

\jieda 如果认为 $y$ 可以精确测定，而所有误差都归结于 $x$，则根据 $\sum_{i=1}^{N}(x_i-\hat{x}_i)^2$ 为最小的原则，可得另一回归方程 $\hat{y} = 70.86\Omega+0.2836x$

\problem{例6-7}

\timu 通过实验测量某量 $x$、$y$ 的结果如下：

\begin{center}
\small
\setlength{\tabcolsep}{5pt}
\begin{tabular}{c|cccccc}
\toprule
$x$ & 2.560 & 2.319 & 2.058 & 1.911 & 1.598 & 0.548\\
\midrule
$y$ & 2.646 & 2.395 & 2.140 & 2.000 & 1.678 & 0.711\\
\bottomrule
\end{tabular}
\end{center}

由重复测量已估计出 $\sigma_x=\sigma_y$，即 $\lambda=1$，试求 $y$ 对 $x$ 的回归直线方程。

\jieda 根据有关公式，计算结果为 $\bar{x}=1.8398, \qquad l_{xx}=2.4495468, \qquad b=0.9747$

$\bar{y}=1.9283, \qquad l_{yy}=2.3273297, \qquad l_{xy}=2.3874952$

$\sum_{i=1}^{N}d_i^2=3.1382\times10^{-4}, \qquad \sigma_y^2=4.023\times10^{-5}, \qquad \sigma_x=\sigma_y=0.0063$

又因为 $b_0=\bar{y}-b\bar{x}=0.1350$

所以所求回归方程为 $\hat{y}=0.1350+0.9747x$

\problem{例6-8}

\timu 用此法说明下列一组数据是否可用 $y=ae^{bx}$ 表示。

\begin{center}
\small
\setlength{\tabcolsep}{5pt}
\begin{tabular}{c|ccccccccc}
\toprule
$x$ & 1 & 2 & 3 & 4 & 5 & 6 & 7 & 8 & 9\\
\midrule
$y$ & 1.78 & 2.24 & 2.74 & 3.74 & 4.45 & 5.31 & 6.92 & 8.85 & 10.97\\
$\log y$ & 0.250 & 0.350 & 0.438 & 0.573 & 0.648 & 0.725 & 0.840 & 0.947 & 1.040\\
\bottomrule
\end{tabular}
\end{center}

\jieda 将$y=ae^{bx}$ 写成线性形式，即 $\log y=\log a+(b\log e)x$

式中，$\log y$ 相当于 $z_1$，$x$ 相当于 $z_2$，$\log a$ 相当于 $A$，$b\log e$ 相当于 $B$。

以 $\log y$ 与 $x$ 画图，选取 $x$ 为 $1,4,6,8$ 的点，可以看出所得图形基本为一直线，故选用函数类型 $y=ae^{bx}$ 是合适的。

\problem{例6-9}

\timu 检验表6-11所示观测数据可用$y=a+be^x$ 表示。

\jieda 具体检验方法如下：第一步，将观测值 $x$ 与 $y$ 画图，得到一条上升曲线；第二步，自曲线上按 $\Delta x$ 为恒定值依次读取 $x$、$y$ 对应值；第三步，列入表中，再依次求出 $\Delta y$、$\log\Delta y$ 以及 $\Delta(\log\Delta y)$。

\begin{center}
\scriptsize
\setlength{\tabcolsep}{5pt}
\begin{tabular}{c|cc|cc|ccc}
\toprule
 & \multicolumn{2}{c|}{观测值} & \multicolumn{2}{c|}{自图上读数值} & \multicolumn{3}{c}{顺序差值}\\
\cline{2-8}
序号 & $x$ & $y$ & $x$ & $y$ & $\Delta y$ & $\log\Delta y$ & $\Delta(\log\Delta y)$\\
\midrule
1 & 0.50 & 17.3 & 0 & 16.6 & --  & --    & --\\
2 & 1.75 & 19.0 & 1 & 17.9 & 1.3 & 0.114 & --\\
3 & 2.75 & 21.0 & 2 & 19.5 & 1.6 & 0.204 & 0.090\\
4 & 3.50 & 22.5 & 3 & 21.5 & 2.0 & 0.301 & 0.097\\
5 & 4.50 & 25.1 & 4 & 23.9 & 2.4 & 0.380 & 0.079\\
6 & 5.25 & 28.0 & 5 & 26.9 & 3.0 & 0.477 & 0.097\\
7 & 6.00 & 30.3 & 6 & 30.6 & 3.7 & 0.568 & 0.091\\
8 & 6.50 & 33.0 & 7 & 35.1 & 4.5 & 0.653 & 0.085\\
9 & 7.50 & 38.0 & 8 & 40.8 & 5.7 & 0.756 & 0.103\\
\bottomrule
\end{tabular}
\end{center}

由表中数据可以看出，$\Delta(\log\Delta y)$ 基本接近常数。因此，所给观测数据可以用 $y=a+be^x$ 表示。

\newpage

\problem{例6-10}

\timu 为了测定椭圆齿轮流量计在介质黏度变化时的误差，先测定 $10$ 号变压器油的黏度 $y$ 与温度 $x$ 的变化曲线，以便试验时测出油温即可知道黏度。通过试验获得如下一组数据：

\begin{center}
\scriptsize
\setlength{\tabcolsep}{4.5pt}
\begin{tabular}{c|ccccccccccccccc}
\toprule
$x/{}^\circ\mathrm{C}$ & 10 & 15 & 20 & 25 & 30 & 35 & 40 & 45 & 50 & 55 & 60 & 65 & 70 & 75 & 80\\
\midrule
$y/{}^\circ\mathrm{E}$ & 4.24 & 3.51 & 2.92 & 2.52 & 2.20 & 2.00 & 1.81 & 1.70 & 1.60 & 1.50 & 1.43 & 1.37 & 1.32 & 1.29 & 1.25\\
\bottomrule
\end{tabular}
\end{center}

试求出黏度与温度之间的经验公式。

\jieda 首先把观测数据点在坐标纸上，并用一条曲线拟合。将该曲线与典型曲线比较，可以看出它很像幂函数 $y=ax^b$ 因此取函数类型为 $y=ax^b$

对等式两边取自然对数，有 $\ln y=\ln a+b\ln x$

令 $y'=\ln(y/{}^\circ\mathrm{E}), \qquad x'=\ln(x/{}^\circ\mathrm{C}), \qquad b_0=\ln a$

则上式变为普通的直线方程 $y'=b_0+bx'$

列表计算如下：

\begin{center}
\scriptsize
\setlength{\tabcolsep}{4pt}
\begin{tabular}{c|c|c|c|c|c|c|c}
\toprule
序号 & $x/{}^\circ\mathrm{C}$ & $y/{}^\circ\mathrm{E}$ & $x'=\ln(x/{}^\circ\mathrm{C})$ & $y'=\ln(y/{}^\circ\mathrm{E})$ & $x'^2$ & $y'^2$ & $x'y'$\\
\midrule
1  & 10 & 4.24 & 2.3026 & 1.4446 & 5.30190 & 2.08676 & 3.32623\\
2  & 15 & 3.51 & 2.7081 & 1.2556 & 7.33354 & 1.57657 & 3.40027\\
3  & 20 & 2.92 & 2.9957 & 1.0716 & 8.97441 & 1.14829 & 3.21018\\
4  & 25 & 2.52 & 3.2189 & 0.9243 & 10.36116 & 0.85425 & 2.97507\\
5  & 30 & 2.20 & 3.4012 & 0.7885 & 11.56814 & 0.62166 & 2.68170\\
6  & 35 & 2.00 & 3.5553 & 0.6931 & 12.64050 & 0.48045 & 2.46438\\
7  & 40 & 1.81 & 3.6889 & 0.5933 & 13.60783 & 0.35204 & 2.18871\\
8  & 45 & 1.70 & 3.8067 & 0.5306 & 14.49068 & 0.28157 & 2.01992\\
9  & 50 & 1.60 & 3.9120 & 0.4700 & 15.30392 & 0.22090 & 1.83867\\
10 & 55 & 1.50 & 4.0073 & 0.4055 & 16.05872 & 0.16440 & 1.62483\\
11 & 60 & 1.43 & 4.0943 & 0.3577 & 16.76366 & 0.12793 & 1.46444\\
12 & 65 & 1.37 & 4.1744 & 0.3148 & 17.42551 & 0.09910 & 1.31414\\
13 & 70 & 1.32 & 4.2485 & 0.2776 & 18.04971 & 0.07708 & 1.17952\\
14 & 75 & 1.29 & 4.3175 & 0.2546 & 18.64070 & 0.06484 & 1.09941\\
15 & 80 & 1.25 & 4.3820 & 0.2231 & 19.20216 & 0.04979 & 0.97782\\
\midrule
$\sum$ & 675 & 30.66 & 54.8134 & 9.6049 & 205.72254 & 8.20563 & 31.76529\\
\bottomrule
\end{tabular}
\end{center}

由表可得 $\sum x'=54.8134, \qquad \bar{x}'=3.6542, \qquad \sum x'^2=205.72254$

$\frac{(\sum x')^2}{N}=200.30080, \qquad l_{x'x'}=\sum x'^2-\frac{(\sum x')^2}{N}=5.42174$

又有 $\sum y'=9.6049, \qquad \bar{y}'=0.6403, \qquad \sum y'^2=8.20563$

$\frac{(\sum y')^2}{N}=6.15034, \qquad l_{y'y'}=\sum y'^2-\frac{(\sum y')^2}{N}=2.05529$

同时 $\sum x'y'=31.76529, \qquad \frac{(\sum x')(\sum y')}{N}=35.09869$

因此 $l_{x'y'} = \sum x'y'-\frac{(\sum x')(\sum y')}{N} = -3.33340$

于是 $b = \frac{l_{x'y'}}{l_{x'x'}} = \frac{-3.33340}{5.42174} = -0.6147$

$b_0 = \bar{y}'-b\bar{x}' = 0.6403+0.6147\times3.6542 = 2.8865$

所以 $\ln(y/{}^\circ\mathrm{E}) = 2.8865-0.6147\ln(x/{}^\circ\mathrm{C})$

即 $\hat{y} = 17.93(x/{}^\circ\mathrm{C})^{-0.6147}{}^\circ\mathrm{E}$

\problem{例6-11}

\timu 对例6-10计算残余平方和 $Q$、残余标准差 $\sigma$ 和相关指数 $R^2$。

\jieda 计算可按下表进行：

\begin{center}
\scriptsize
\setlength{\tabcolsep}{5pt}
\begin{tabular}{c|c|c|c|c|c}
\toprule
序号 & $y/{}^\circ\mathrm{E}$ & $y^2/{}^\circ\mathrm{E}^2$ & $\hat{y}/{}^\circ\mathrm{E}$ & $(y-\hat{y})/{}^\circ\mathrm{E}$ & $(y-\hat{y})^2/{}^\circ\mathrm{E}^2$\\
\midrule
1  & 4.24 & 17.9776 & 4.355 & $-0.115$ & 0.013225\\
2  & 3.51 & 12.3201 & 3.394 & 0.116 & 0.013456\\
3  & 2.92 & 8.5264  & 2.844 & 0.076 & 0.005776\\
4  & 2.52 & 6.3504  & 2.479 & 0.041 & 0.001681\\
5  & 2.20 & 4.8400  & 2.216 & $-0.016$ & 0.000256\\
6  & 2.00 & 4.0000  & 2.016 & $-0.016$ & 0.000256\\
7  & 1.81 & 3.2761  & 1.857 & $-0.047$ & 0.002209\\
8  & 1.70 & 2.8900  & 1.727 & $-0.027$ & 0.000729\\
9  & 1.60 & 2.5600  & 1.619 & $-0.019$ & 0.000361\\
10 & 1.50 & 2.2500  & 1.527 & $-0.027$ & 0.000729\\
11 & 1.43 & 2.0449  & 1.447 & $-0.017$ & 0.000289\\
12 & 1.37 & 1.8769  & 1.378 & $-0.008$ & 0.000064\\
13 & 1.32 & 1.7427  & 1.316 & 0.004 & 0.000016\\
14 & 1.29 & 1.6641  & 1.262 & 0.028 & 0.000784\\
15 & 1.25 & 1.5625  & 1.213 & 0.037 & 0.001369\\
\midrule
$\sum$ & 30.66 & 73.8814 & -- & 0.010 & 0.041200\\
\bottomrule
\end{tabular}
\end{center}

因此 $Q = \sum_{i=1}^{N}(y_i-\hat{y}_i)^2 = 0.0412\,{}^\circ\mathrm{E}^2$

残余标准差为 $\sigma = \sqrt{\frac{Q}{N-2}} = \sqrt{\frac{0.0412}{15-2}}\,{}^\circ\mathrm{E} = 0.056{}^\circ\mathrm{E}$

总离差平方和为
\begin{equation*}
    \sum_{i=1}^{N}(y_i-\bar{y})^2 = \sum_{i=1}^{N}y_i^2 - \frac{\left(\displaystyle\sum_{i=1}^{N}y_i\right)^2}{N} = 11.212\,{}^\circ\mathrm{E}^2
\end{equation*}

所以相关指数为
\begin{equation*}
    R^2 = 1- \frac{Q}{\displaystyle\sum_{i=1}^{N}(y_i-\bar{y})^2} = 0.9963
\end{equation*}

\newpage

\problem{例6-12}

\timu 根据经验知道某变量 $y$ 受变量 $x_1$、$x_2$ 影响，通过试验获得下表中的一批数据，试建立 $y$ 对 $x_1$、$x_2$ 的线性回归方程。

\begin{center}
\scriptsize
\setlength{\tabcolsep}{4pt}

\begin{tabular}{c|ccccccccccccccc}
\toprule
序号 & 1 & 2 & 3 & 4 & 5 & 6 & 7 & 8 & 9 & 10 & 11 & 12 & 13 & 14 & 15\\
\midrule
$x_1$ & 15.58 & 10.68 & 15.62 & 15.78 & 13.22 & 16.44 & 11.40 & 16.17 & 14.03 & 15.67 & 12.74 & 11.73 & 14.84 & 13.73 & 15.12\\
$x_2$ & 1.95 & 1.37 & 2.39 & 1.14 & 1.85 & 1.32 & 2.05 & 1.11 & 1.47 & 1.38 & 1.35 & 1.33 & 1.09 & 1.27 & 1.78\\
$y$   & 1.34 & 1.27 & 1.56 & 1.48 & 1.40 & 1.82 & 0.85 & 1.40 & 1.15 & 1.89 & 0.87 & 1.53 & 1.25 & 2.47 & 1.83\\
\bottomrule
\end{tabular}

\vspace{0.5em}

\begin{tabular}{c|cccccccccccccc}
\toprule
序号 & 16 & 17 & 18 & 19 & 20 & 21 & 22 & 23 & 24 & 25 & 26 & 27 & 28 & 29\\
\midrule
$x_1$ & 17.88 & 13.38 & 14.21 & 16.80 & 16.38 & 10.81 & 17.26 & 14.92 & 18.14 & 18.15 & 10.31 & 11.40 & 12.57 & 17.61\\
$x_2$ & 2.52 & 1.43 & 2.27 & 1.41 & 1.78 & 1.32 & 1.31 & 1.42 & 2.13 & 1.20 & 0.98 & 1.27 & 0.87 & 1.21\\
$y$   & 2.41 & 1.69 & 1.59 & 1.19 & 2.44 & 1.35 & 1.57 & 1.64 & 1.64 & 2.34 & 0.65 & 1.19 & 2.06 & 1.57\\
\bottomrule
\end{tabular}

\end{center}

\jieda 设欲求的回归方程形式为 $\hat{y} = \mu_0+b_1(x_1-\bar{x}_1)+b_2(x_2-\bar{x}_2)$

由数据计算得 $N=29, \qquad \sum x_1=422.57, \qquad \sum x_2=43.97, \qquad \sum y=45.44$

$\bar{x}_1=14.571, \qquad \bar{x}_2=1.516, \qquad \bar{y}=1.567$

同时有 $\sum x_1^2=6317.3183, \qquad \sum x_1x_2=647.9873, \qquad \sum x_2^2=71.8675$

$\sum x_1y=678.2800, \qquad \sum x_2y=69.7515, \qquad \sum y^2=77.2434$

因此 $l_{11} = \sum x_1^2-\frac{1}{N}\left(\sum x_1\right)^2 = 159.8906$

$l_{12} = \sum x_1x_2-\frac{1}{N}\left(\sum x_1\right)\left(\sum x_2\right) = 7.2838$

$l_{22} = \sum x_2^2-\frac{1}{N}\left(\sum x_2\right)^2 = 5.1999$

$l_{1y} = \sum x_1y-\frac{1}{N}\left(\sum x_1\right)\left(\sum y\right) = 16.1565$

$l_{2y} = \sum x_2y-\frac{1}{N}\left(\sum x_2\right)\left(\sum y\right) = 0.8551$

$l_{yy} = \sum y^2-\frac{1}{N}\left(\sum y\right)^2 = 6.0436$

其矩阵 $L$ 为
\begin{equation*}
L=
    \begin{pmatrix}
        l_{11} & l_{12}\\
        l_{21} & l_{22}
    \end{pmatrix}
    =
    \begin{pmatrix}
        159.8906 & 7.2838\\
        7.2838 & 5.1999
    \end{pmatrix}
\end{equation*}

其逆矩阵为
\begin{equation*}
L^{-1}
    =
    \begin{pmatrix}
        0.006681 & -0.009358\\
        -0.009358 & 0.205419
    \end{pmatrix}
\end{equation*}

于是 $\mu_0=\bar{y}=1.567$

\begin{equation*}
\begin{pmatrix}
        b_1\\
        b_2
    \end{pmatrix}
    =
    \begin{pmatrix}
        0.006681 & -0.009358\\
        -0.009358 & 0.205419
    \end{pmatrix}
    \begin{pmatrix}
        16.1565\\
        0.8551
    \end{pmatrix}
\end{equation*}

即 $b_1=0.006681\times16.1565-0.009358\times0.8551=0.0999$

$b_2=-0.009358\times16.1565+0.205419\times0.8551=0.0245$

所以所求回归方程为 $\hat{y} = 1.567+0.0999(x_1-14.571)+0.0245(x_2-1.516)$

\problem{例6-13}

\timu 今测得一组小学生的身高 $x_1$、年龄 $x_2$ 和体重 $y$ 的数据如下表，试求体重、身高与年龄的回归关系并进行检验。

\begin{center}
\scriptsize
\setlength{\tabcolsep}{4pt}
\begin{tabular}{c|cccccccccccc}
\toprule
$x_1/\mathrm{cm}$ & 147 & 149 & 139 & 152 & 141 & 140 & 145 & 138 & 142 & 132 & 151 & 147\\
\midrule
$x_2/\text{岁}$ & 9 & 11 & 7 & 12 & 9 & 8 & 11 & 10 & 11 & 7 & 13 & 10\\
\midrule
$y/\mathrm{kg}$ & 34 & 41 & 23 & 37 & 25 & 28 & 47 & 27 & 26 & 21 & 46 & 38\\
\bottomrule
\end{tabular}
\end{center}

\jieda 设欲求的回归方程形式为 $\hat{y} = \mu_0+b_1(x_1-\bar{x}_1)+b_2(x_2-\bar{x}_2)$

由数据计算得 $N=12, \qquad \sum x_1=1723, \qquad \sum x_2=118, \qquad \sum y=393$

$\bar{x}_1=143.5833\approx143.58, \qquad \bar{x}_2=9.8333\approx9.83, \qquad \bar{y}=32.75$

并有 $\sum x_1^2=247783, \qquad \sum x_2^2=1200, \qquad \sum x_1x_2=17042$

$\sum y^2=13759, \qquad \sum x_1y=56910, \qquad \sum x_2y=4009$

于是 $l_{11} = \sum x_1^2-\frac{1}{N}\left(\sum x_1\right)^2 = 388.92$

$l_{22} = \sum x_2^2-\frac{1}{N}\left(\sum x_2\right)^2 = 39.67$

$l_{12} = \sum x_1x_2-\frac{1}{N}\left(\sum x_1\right)\left(\sum x_2\right) = 99.17$

$l_{yy} = \sum y^2-\frac{1}{N}\left(\sum y\right)^2 = 888.25$

$l_{1y} = \sum x_1y-\frac{1}{N}\left(\sum x_1\right)\left(\sum y\right) = 481.75$

$l_{2y} = \sum x_2y-\frac{1}{N}\left(\sum x_2\right)\left(\sum y\right) = 144.5$

其矩阵 $L$ 为
\begin{equation*}
L=
    \begin{pmatrix}
        388.92 & 99.17\\
        99.17 & 39.67
    \end{pmatrix}
\end{equation*}

其逆矩阵为
\begin{equation*}
L^{-1}
    =
    \begin{pmatrix}
        0.0071 & -0.0177\\
        -0.0177 & 0.0695
    \end{pmatrix}
\end{equation*}

根据公式得 $\mu_0=\bar{y}=32.75$

\begin{equation*}
\begin{pmatrix}
        b_1\\
        b_2
    \end{pmatrix}
    =
    L^{-1}
    \begin{pmatrix}
        l_{1y}\\
        l_{2y}
    \end{pmatrix}
    =
    \begin{pmatrix}
        0.0071 & -0.0177\\
        -0.0177 & 0.0695
    \end{pmatrix}
    \begin{pmatrix}
        481.75\\
        144.5
    \end{pmatrix}
\end{equation*}

即 $b_1=0.0071\times481.75-0.0177\times144.5=0.8546\approx0.85$

$b_2=-0.0177\times481.75+0.0695\times144.5=1.5065\approx1.51$

所求回归方程为 $\hat{y} = 32.75+0.85(x_1-143.58)+1.51(x_2-9.83)$

\newpage

\problem{例6-14}

\timu 对例6-12中的回归进行方差分析。

\jieda 方差分析表如下：

\begin{center}
\small
\setlength{\tabcolsep}{5pt}
\begin{tabular}{c|c|c|c|c|c}
\toprule
来源 & 平方和 & 自由度 & 方差 & $F$ & 显著性\\
\midrule
回归 & 1.6350 & 2 & 0.8175 & 4.82 & $\alpha=0.05$\\
残余 & 4.4086 & 26 & 0.1696 & -- & --\\
总计 & 6.0436 & 28 & -- & -- & --\\
\bottomrule
\end{tabular}
\end{center}

总离差平方和为 $S=l_{yy}=6.0436$，自由度为 $\nu_S=29-1=28$；回归平方和为 $U=b_1l_{1y}+b_2l_{2y}=1.6350$，自由度为 $\nu_U=2$；残余平方和为 $Q=S-U=4.4086$，自由度为 $\nu_Q=28-2=26$。

由于 $F=4.82>F_{0.05}(2,26)=3.37$，因此回归方程 $\hat{y}=0.074+0.0999x_1+0.0245x_2$ 在 $\alpha=0.05$ 水平上显著。

残余标准差为 $\sigma=\sqrt{0.1696}=0.41$，于是 $2\sigma=0.82$。用该回归方程进行预报时，$95\%$ 的误差不会超过 $0.82$。

\newpage

\problem{例6-15}

\timu 对例6-13中的回归进行方差分析。

\jieda 方差分析表如下：

\begin{center}
\small
\setlength{\tabcolsep}{5pt}
\begin{tabular}{c|c|c|c|c|c}
\toprule
来源 & 平方和 & 自由度 & 方差 & $F$ & 显著性\\
\midrule
回归 & 627.6825 & 2 & 313.8413 & 10.84 & $F_{0.01}=8.02$\\
残余 & 260.5675 & 9 & 28.9519 & -- & --\\
总计 & 888.25 & 11 & -- & -- & --\\
\bottomrule
\end{tabular}
\end{center}

总离差平方和为 $S=l_{yy}=888.25$，自由度为 $\nu_S=12-1=11$；回归平方和为 $U=b_1l_{1y}+b_2l_{2y}=627.6825$，自由度为 $\nu_U=2$；残余平方和为 $Q=l_{yy}-U=260.5675$，自由度为 $\nu_Q=11-2=9$。

由于 $F=10.84>F_{0.01}(2,9)=8.02$，因此回归方程 $\hat{y}=0.85x_1+1.51x_2-104.14$ 在 $\alpha=0.01$ 水平上高度显著。

残余标准差为 $\sigma=5.38\mathrm{kg}$，于是 $2\sigma=10.76\mathrm{kg}$。用该回归方程进行预报时，$99\%$ 的误差不会超过 $10.76\mathrm{kg}$。

\problem{例6-16}

\timu 分析例6-12中两个自变量在回归中所起的作用。

\jieda 按有关公式计算：
\begin{equation*}
    P_1=\frac{b_1^2}{c_{11}}=\frac{0.0999^2}{0.006681}=1.4938,\quad P_2=\frac{b_2^2}{c_{22}}=\frac{0.0245^2}{0.205419}=0.0029
\end{equation*}

又因为例6-14中残余方差为 $\sigma^2=0.1696$

所以
\begin{equation*}
    F_1=\frac{P_1}{\sigma^2} = \frac{1.4938}{0.1696} = 8.81,\quad F_2=\frac{P_2}{\sigma^2} = \frac{0.0029}{0.1696} = 0.017
\end{equation*}

查 $F$ 分布表，得 $F_{0.1}(1,26)=2.91, \qquad F_{0.01}(1,26)=7.72$

因为 $F_1>F_{0.01}(1,26)$ 所以 $P_1$ 在 $\alpha=0.01$ 水平上显著。

又因为 $F_2<F_{0.1}(1,26)$ 所以 $P_2$ 不显著。这说明 $x_1$ 是影响 $y$ 的一个主要因素，而 $x_2$ 对影响很小。

因此，可以在回归方程 $\hat{y}=1.567+0.0999(x_1-14.571)+0.0245(x_2-1.516)$ 中把自变量 $x_2$ 删除。此时，$y$ 对 $x_1$ 的回归系数 $b_1'$ 可按 $b_1' = b_1-\frac{c_{12}}{c_{22}}b_2$ 求得，即 $b_1' = 0.0999- \frac{-0.009358}{0.205419}\times0.0245 = 0.1010$

从而得到新的回归方程 $\hat{y} = 1.567+0.1010(x_1-14.571)$

\end{document}
